我们现在假设 \(\varphi_{\lambda}\) 在无穷远处不趋于 0，并记 \(\theta = 1 / \lambda\)。对每个实数 \(x\)，记 \(\langle x\rangle\) 为 \([0,1[\) 中满足 \(x - \langle x\rangle \in \mathbf{Z}\) 的唯一实数。

\begin{enumerate}
\def\labelenumi{\alph{enumi})}
\setcounter{enumi}{7}
\item
  存在实数 \(u \neq 0\)，使得通项为 \(\sin(u\theta^{n})^{2}\) 的级数收敛；且通项为 \(\langle u\theta^{n}\rangle\) 的级数收敛。
\item
  对每个 \(n \in N\)，设 \(c_{n}\) 为满足 \(|u\theta^{n} - c_{n}| \leqslant \frac{1}{2}\) 的整数。形式级数 \(g(z) = \sum c_{n}z^{n}\) 是一个有理函数。（利用 A，IV，第 85 页，练习 1 中的判别法以及 Hadamard 不等式（EVT，V，第 37 页，推论 3），证明所出现的那些 Hankel 行列式趋于 0；再注意到这些行列式都是整数，由此得出结论。）
\item
  存在整系数多项式 \(f_1\) 与 \(f_2\)，使得 \(g = f_1 / f_2\) 且 \(f_2(0) = 1\)。
\item
  数 \(\theta\) 是一个 Pisot 数（“Salem 定理”）。（注意 \(f_{2}(\lambda)=0\)，并且 \(f_{2}\) 的其余根的模均 \textgreater1。）
\end{enumerate}

\begin{enumerate}
\def\labelenumi{\arabic{enumi})}
\setcounter{enumi}{58}
\tightlist
\item
  设 \(n \geqslant 1\) 为整数，并为 \(R^{n}\) 配备 Euclid 范数。我们用 \(\mu\) 表示 \(R^{n}\) 上的 Lebesgue 测度。
\end{enumerate}

设 r \textgreater{} 0。\(R^{n}\) 中半径为 r 的球堆积是指由 \(R^{n}\) 的两两不相交的子集组成的一个集合 P，使得每个 \(S \in P\) 都是半径为 r 的闭球。我们称 P 是周期的，如果存在一个格 \(\Lambda \subset R^{n}\)，使得 \(S + h \in P\) 对所有 \(S \in P\) 以及每个 \(h \in \Lambda\) 都成立。

设 \(A _{P}\) 为 P 的诸元素之并。它是一个可测集。我们称 P 的密度为实数

\[
\Delta (\mathcal {P}) = \limsup _ {\mathrm{R} \to + \infty} \frac {\mu (\mathrm{A} _ {\mathcal {P}} \cap \mathrm{B} _ {\mathrm{R}})}{\mu (\mathrm{B} _ {\mathrm{R}})}
\]

其中 \(B_{R}\) 是 \(R^{n}\) 中半径为 R 的闭球。我们记 \(\Delta_{n} = \sup_{\mathcal{P}} \Delta(\mathcal{P})\)。a) 我们有

\[
\Delta_ {n} = \sup _ {\mathcal {P} \text {   periodic }} \Delta (\mathcal {P}).
\]

\begin{enumerate}
\def\labelenumi{\alph{enumi})}
\setcounter{enumi}{1}
\tightlist
\item
  设 P 为一个周期的球堆积，它在格 \(\Lambda\) 下不变。设 \(C \subset R^{n}\) 为诸球 \(S \in P\) 的中心模 \(\Lambda\) 的一个代表集。我们有
\end{enumerate}

\[
\Delta (\mathcal {P}) = \mu (\mathrm{B} _ {1}) \frac {r ^ {n} \operatorname{Card} (\mathrm{C})}{\mathrm{V} (\Lambda)}
\]

其中 V(Λ) 是 Λ 的余体积。

\begin{enumerate}
\def\labelenumi{\alph{enumi})}
\setcounter{enumi}{2}
\tightlist
\item
  设 \(a > 0\) 为一个实数。设 \(f: \mathbf{R}^{n} \to \mathbf{R}\) 为一个 Schwartz 函数，使得
\end{enumerate}

\begin{enumerate}
\def\labelenumi{(\roman{enumi})}
\item
  我们有 \(f(x) \leqslant 0 \text{ 若 } \|x\| \geqslant 2a\)；
\item
  \(f\) 的 Fourier 变换取实值且为正，并满足 \(\widehat{f}(0)>0\)。
\end{enumerate}

则我们有

\[
\Delta_ {n} \leqslant \mu (\mathrm{B} _ {1}) a ^ {n} \frac {f (0)}{\widehat {f} (0)}.
\]

（只需考虑半径为 \(a\) 的周期球堆积；设 \(\mathcal{P}\) 为其中之一，并设 \(\Lambda\) 与 \(C \subset \mathbf{R}^n\) 如上一问题中那样；利用 Poisson 公式证明

\[
\sum_ {(c, d) \in \mathrm{C} ^ {2}} \sum_ {n \in \Lambda} f (d - c + n) = \frac {1}{\mathrm{V} (\Lambda)} \sum_ {h \in \Lambda^ {*}} \left| \sum_ {c \in \mathrm{C}} e ^ {2 i \pi c \cdot h} \right| ^ {2} \widehat {f} (h),
\]

其中 \(\Lambda^{*}\) 是相伴的格，然后得出结论。）

\begin{enumerate}
\def\labelenumi{\arabic{enumi})}
\setcounter{enumi}{59}
\tightlist
\item
  \begin{enumerate}
  \def\labelenumii{\alph{enumii})}
  \tightlist
  \item
    设 \(\mathbf{H} \subset \mathbf{C}\) 为由复数 \(z\)（满足 \(\mathcal{I}(z) > 0\)）组成的开集。函数
  \end{enumerate}
\end{enumerate}

\[
\Theta (z) = \sum_ {n \in \mathbf {Z}} e ^ {2 i \pi n ^ {2} z}
\]

在 H 上全纯，且它满足 \(\Theta(z) = (z/i)^{-1/2}\Theta(-1/z)\)（对 \(z\in\mathrm{H}\)），其中 \(z\mapsto(z/i)^{1/2}\) 是 H 上延拓函数 \(iy\mapsto\sqrt{y}\) 的唯一的全纯函数。（利用练习 16，b)。）

\begin{enumerate}
\def\labelenumi{\alph{enumi})}
\setcounter{enumi}{1}
\tightlist
\item
  设 m 与 n 为严格正的整数。我们有
\end{enumerate}

\[
\frac {1}{\sqrt {m}} \sum_ {k = 0} ^ {m - 1} e ^ {2 i \pi k ^ {2} n / m} = \frac {e ^ {i \pi / 4}}{\sqrt {2 n}} \sum_ {k = 0} ^ {2 n - 1} e ^ {- 2 i \pi k ^ {2} m / (2 n)}.
\]

（在上一问题的关系式中取 \(z = -2in / m + iy\)，并令 \(y \to 0\)。）

\begin{enumerate}
\def\labelenumi{\alph{enumi})}
\setcounter{enumi}{2}
\tightlist
\item
  对每个整数 \(m \geqslant 1\) 以及每个 \(n \in Z\)，置
\end{enumerate}

\[
\tau (n, m) = \sum_ {k = 0} ^ {m - 1} e ^ {2 i \pi k ^ {2} n / m}.
\]

对所有严格正的整数 \(m_{1}\) 与 \(m_{2}\)，我们有

\[
\tau (1, m _ {1} m _ {2}) = \tau (m _ {2}, m _ {1}) \tau (m _ {1}, m _ {2}).
\]

\begin{enumerate}
\def\labelenumi{\alph{enumi})}
\setcounter{enumi}{3}
\item
  若 m 为奇数，我们有 \(\tau(1,m)=i^{((m-1)/2)^{2}}\sqrt{m}\)。（利用 a)。
\item
  若 p 是一个奇素数，则 \(\tau(n,p)=\ell_{p}(n)\tau(1,p)\)，其中 \(\ell_{p}\) 是 \((\mathbf{Z}/p\mathbf{Z})^{*}\) 的唯一的 2 阶特征标，被延拓到 \(Z/pZ\)（通过置 \(\ell_{p}(0)=0\)）。
\item
  设 \(p \neq q\) 为两个奇素数。我们有
\end{enumerate}

\[
\ell_ {p} (q) \ell_ {q} (p) = \frac {\tau (1 , p q)}{\tau (1 , p) \tau (1 , q)}.
\]

\begin{enumerate}
\def\labelenumi{\alph{enumi})}
\setcounter{enumi}{6}
\tightlist
\item
  由此给出二次互反律的一个新证明（参看 A，V，第 156 页，练习 23）：对所有满足 \(p \neq q\) 的奇素数，有
\end{enumerate}

\[
\ell_ {p} (q) \ell_ {q} (p) = (- 1) ^ {(p - 1) (q - 1) / 4}.
\]

\begin{enumerate}
\def\labelenumi{\arabic{enumi})}
\setcounter{enumi}{60}
\tightlist
\item
  设 \(r \geqslant 1\) 为整数。设 \(A = (a_{i,j})\) 为一个 \((r, r)\) 型的整系数对称正定矩阵，并设 \(N \geqslant 1\) 为使得矩阵 \(A^* = NA^{-1}\) 也有整系数的整数。我们用 Q（相应地，\(\widetilde{Q}\)）表示与 A（相应地，与 \(A^{-1}\)）相伴的整二次型，也就是说，\(Q(x) = {}^t xAx\) 对所有 \(x \in \mathbf{Z}^r\) 成立。我们用 \((a, b) \mapsto \langle a|b \rangle_A\) 表示与 A 相伴的双线性型，于是 \(\langle a|b \rangle_A = {}^t aAb\)。
\end{enumerate}

我们用 H 表示由复数 z（满足 \(\mathcal{I}(z)>0\)）组成的集合。它是 C 的一个开子集。我们用 log: \(H\to C\) 表示对数的主值分支在 H 上的限制（FVR，III，第 10 页）。设 \(k\in N\) 为奇整数。对 \(z\in H\)，我们记 \(z^{k}=\exp(k\log(z))\)。

\begin{enumerate}
\def\labelenumi{\alph{enumi})}
\tightlist
\item
  群 \(\mathrm{SL}_2(\mathbf{Z})\) 按下式传递地作用于 \(\mathbf{H}\)：
\end{enumerate}

\[
\left( \begin{array}{c c} a & b \\ c & d \end{array} \right) \cdot z = \frac {a z + b}{c z + d}, \qquad \left( \begin{array}{c c} a & b \\ c & d \end{array} \right) \in \mathrm{SL} _ {2} (\mathbf {Z})
\]

  （参看 TA，I，第 150 页，练习 5）。

\begin{enumerate}
\def\labelenumi{\alph{enumi})}
\setcounter{enumi}{1}
\tightlist
\item
  对每个复数 \(z \in \mathbf{H}\) 以及每个 \(x \in \mathbf{Z}^r\)，级数
\end{enumerate}

\[
\sum_{\substack{m\in \mathbf{Z}^{r}\\ m\equiv x\bmod \mathrm{N}}}\exp (i\pi \mathrm{Q}(m)z)
\]

绝对收敛且在紧集上一致收敛。

我们用 \(\theta_{\mathrm{A}}(z;x)\) 或 \(\theta_{\mathrm{Q}}(z;x)\) 表示这个级数的和（“与 Q 相伴的 θ 函数”）。它是 \(\mathbf{H}\) 上的一个全纯函数。

\begin{enumerate}
\def\labelenumi{\alph{enumi})}
\setcounter{enumi}{2}
\tightlist
\item
  对每个 \(y \in \mathbf{Z}^r\)，有
\end{enumerate}

\[
\sum_ {m \in {\bf Z} ^ {r}} \exp (i \pi Q (m + y) z) = \frac {1}{\sqrt {\det (A)}} \Bigl (\frac {i}{z} \Bigr) ^ {r / 2} \sum_ {m \in {\bf Z} ^ {r}} \exp \Bigl (- i \pi \frac {\widetilde {Q} (m)}{z} + m \cdot y \Bigr),
\]

（应用 Poisson 公式与练习 1。）

\begin{enumerate}
\def\labelenumi{\alph{enumi})}
\setcounter{enumi}{3}
\tightlist
\item
  我们用 H 表示 \((\mathbf{Z}/\mathrm{NZ})^{r}\) 的由元素 \(\alpha\in(\mathbf{Z}/\mathrm{NZ})^{r}\)（满足 \(A\alpha=0\)，在 Z/NZ 中）组成的子群。群 H 的基数为 det(A)。H 的每个特征标都具有如下形式
\end{enumerate}

\[
\psi_ {\beta} (\alpha) = \exp \Bigl (2 i \pi \frac {\langle \alpha | \beta \rangle_ {\mathrm{A}}}{\mathrm{N} ^ {2}} \Bigr),
\]

其中 \(\beta\in\mathrm{H}\)。

\begin{enumerate}
\def\labelenumi{\alph{enumi})}
\setcounter{enumi}{4}
\item
  对每个 \(\alpha \in \mathrm{H}\) 与 \(z\in \mathbf{H}\)，有 \(\theta_{\mathrm{Q}}(z + 2;\alpha) = \exp \bigl (\frac{2i\pi\mathrm{Q}(\alpha)}{\mathrm{N}^2}\bigr)\theta_{\mathrm{Q}}(z;\alpha)\)。
\item
  对每个 \(\alpha \in \mathrm{H}\) 与 \(z\in \mathbf{H}\)，有
\end{enumerate}

\[
\theta_ {\mathrm{Q}} \Bigl (- \frac {1}{z}; \alpha \Bigr) = \frac {1}{\sqrt {\det (\mathrm{A})}} (- i z) ^ {r / 2} \sum_ {\beta \in \mathrm{H}} \psi_ {\beta} (\alpha) \theta_ {\mathrm{Q}} (z; \beta).
\]

\(g)\) 设 \(g = \begin{pmatrix} a & b \\ c & d \end{pmatrix} \in \mathrm{SL}_2(\mathbf{Z})\) 满足 \(c \neq 0\)、\(d > 0\) 且 \(b \equiv c \equiv 0 \mod 2\)。对 \(\alpha \in \mathrm{H}\) 与 \(z \in \mathbf{H}\)，有

\[
\theta_ {\mathrm{Q}} (g \cdot z; \alpha) = \frac {1}{d ^ {r / 2} \det (\mathrm{A})} (c z + d) ^ {r / 2} \sum_ {\beta \in \mathrm{H}} \Phi (\alpha , \beta) \theta_ {\mathrm{Q}} (z; \beta),
\]

其中

\[
\Phi (\alpha ,\beta) = \sum_{\gamma \in \mathrm{H}}\psi_{\gamma}(\alpha)\Bigl (\sum_{\substack{\delta \text{mod.} d\mathrm{N}\\ \delta \equiv \alpha \text{mod.}\mathrm{N}}}\exp \Bigl (\frac{i\pi}{d\mathrm{N}^{2}} (b\mathrm{Q}(\delta) + 2\langle \gamma |\delta \rangle_{\mathrm{A}} - c\mathrm{Q}(\gamma))\Bigr)\Bigr)
\]

（先考虑 \(\theta_{\mathrm{Q}}(\widetilde{g} \cdot z; \alpha)\)，其中 \(\widetilde{g} = g\begin{pmatrix} 0 & -1 \\ 1 & 0 \end{pmatrix}\)。注意公式 \(d\widetilde{g} \cdot z = b - (dz - c)^{-1}\) 并应用上一问题；注意 \(\Phi(\alpha, \beta)\) 只依赖于 \(\beta\) 模 N 的值。）

\begin{enumerate}
\def\labelenumi{\alph{enumi})}
\setcounter{enumi}{7}
\tightlist
\item
  再设 \(c \equiv 0 \mod. 2\mathrm{N}\)。则
\end{enumerate}

\[
\theta_ {\mathrm{Q}} (g \cdot z; x) = \frac {\tau}{d ^ {r / 2}} (c z + d) ^ {k} \theta_ {\mathrm{Q}} (z; a \alpha),
\]

其中

\[
\tau = \sum_{\substack{\delta \text{mod.} d\mathrm{N}\\ \delta \equiv x\text{mod.}\mathrm{N}}}\exp \Bigl (\frac{i\pi b\mathrm{Q}(\delta)}{d\mathrm{N}^{2}}\Bigr).
\]

\begin{enumerate}
\def\labelenumi{\roman{enumi})}
\tightlist
\item
  再设 \(d\) 为奇数。我们有 \(\tau = \exp \left(\frac{i\pi abQ(\alpha)}{N^2}\right)\tau_Q(c;d)\)，其中
\end{enumerate}

\[
\tau_ {\mathrm{Q}} (c, d) = \sum_ {x \in (\mathbf {Z} / d \mathbf {Z}) ^ {r}} e \left(- \frac {i \pi c \mathrm{Q} (x)}{d}\right)
\]

（“Gauss 和”。）

\begin{enumerate}
\def\labelenumi{\alph{enumi})}
\setcounter{enumi}{9}
\tightlist
\item
  设 \(d = p_{1} \cdots p_{k}\)，其中 \(p_{1}, \ldots, p_{k}\) 为素数。我们置
\end{enumerate}

\[
\left(\frac {n}{d}\right) = \prod_ {i = 1} ^ {k} \ell_ {p _ {i}} (n)
\]

其中，对每个奇素数 \(p\)，我们用 \(\ell_p\) 表示 \((\mathbf{Z}/p\mathbf{Z})^*\) 的唯一的 2 阶特征标，它被延拓到 \(\mathbf{Z}/p\mathbf{Z}\)（通过置 \(\ell_p(0) = 0\)；“Jacobi 符号”）。若 \(d\) 与 \(2c\det(\mathrm{A})\) 互素，则

\[
\tau_ {\mathrm{Q}} (c, d) = \Big (\frac {\det (\mathrm{A})}{d} \Big) \Big (\overline {{\varepsilon}} _ {d} \Big (\frac {2 c}{d} \Big) \sqrt {d} \Big) ^ {r},
\]

其中

\[
\varepsilon_ {d} = \left\{ \begin{array}{l l} 1 & \text { if } d \equiv 1 \bmod . 4 \\ i & \text { if } d \equiv - 1 \bmod . 4. \end{array} \right.
\]

（将二次型对角化并利用上一练习。）

\(k)\) 若 \(g\) 模 4N 与单位矩阵同余，则

\[
\theta_ {\mathrm{Q}} (g \cdot z; \alpha) = \left(\frac {2 c}{d}\right) ^ {r} (c z + d) ^ {k} \theta_ {\mathrm{Q}} (z; \alpha).
\]

\(^{*}\) 这意味着映射 \(z \mapsto \theta_{q}(z; \alpha)\) 是一个权为 r/2、水平为 4N 的全纯模形式。（FM，编写中。）\(^{*}\)

\begin{enumerate}
\def\labelenumi{\arabic{enumi})}
\setcounter{enumi}{61}
\tightlist
\item
  我们把 \(\mathbf{T}\) 与区间 \([-1/2, 1/2]\) 等同。设 \(f\) 是 \(\mathbf{R}/\mathbf{Z}\) 上由 \(f(x) = s(x)\)（符号函数）定义的函数。设 \(s_n\) 为对称部分和
\end{enumerate}

\[
s _ {n} (x) = \sum_ {h = - n} ^ {n} \widehat {f} (h) e ^ {2 i \pi h x}
\]

，即 f 的 Fourier 级数的对称部分和。

\begin{enumerate}
\def\labelenumi{\alph{enumi})}
\item
  对每个 \(x \neq 0\)（在 \(\mathbf{T}\) 中），有 \(s_n(x) \to f(x)\)（当 \(n \to +\infty\)）。
\item
  设 \(n \geqslant 1\) 且 \(0 < x \leqslant 1/2\)。我们有
\end{enumerate}

\[
s _ {2 n} (x) = \frac {\sin (2 \pi n x)}{2 \sin (\pi x)}.
\]

\begin{enumerate}
\def\labelenumi{\alph{enumi})}
\setcounter{enumi}{2}
\tightlist
\item
  我们有
\end{enumerate}

\[
\lim _ {n \rightarrow + \infty} s _ {2 n} \left(\frac {1}{4 n}\right) = 2 \int_ {0} ^ {1 / 2} \frac {\sin (2 \pi x)}{\pi x} d x,
\]

并且这个极限属于区间 \(]1,2[\)（“Gibbs 现象”）。

\begin{enumerate}
\def\labelenumi{\arabic{enumi})}
\setcounter{enumi}{62}
\tightlist
\item
  设 G 为 \(T^{N}\) 的由元素 \(x = (x_{n})\)（满足：除至多有限多个例外，\(2x_{n} = 0\) 对所有 \(n \in N\) 成立）组成的子群。
\end{enumerate}

\begin{enumerate}
\def\labelenumi{\alph{enumi})}
\item
  对每个有限集 \(\Lambda \subset \mathbf{N}\)，设 \(\mathrm{U}_{\Lambda} \subset \mathrm{G}\) 为由满足下列条件的 \((x_{n})\) 组成的集合：\(x_{n} = 1\)（对 \(n \in \Lambda\)）且 \(2x_{n} = 0\)（对 \(n \notin \Lambda\)）。G 上存在一个拓扑群结构，使得集合族 \(U_{\Lambda}\) 构成 e 的开邻域基。
\item
  拓扑群 G 是局部紧的。
\item
  群 G 不是可除的。
\end{enumerate}

\(d)\) 对偶群 \(\widehat{\mathrm{G}}\) 是无挠的。

\begin{enumerate}
\def\labelenumi{\arabic{enumi})}
\setcounter{enumi}{63}
\tightlist
\item
  设 \(\mu\) 为 \(\mathbf{T}\) 上的一个测度。
\end{enumerate}

\begin{enumerate}
\def\labelenumi{\alph{enumi})}
\tightlist
\item
  对每个 \(x \in \mathbf{T}\)，有
\end{enumerate}

\[
\mu (\{x \}) = \lim _ {\mathrm{H} \rightarrow + \infty} \frac {1}{2 \mathrm{H} + 1} \sum_ {| h | \leqslant \mathrm{H}} \widehat {\mu} (h) e ^ {2 i \pi h x}.
\]

\begin{enumerate}
\def\labelenumi{\alph{enumi})}
\setcounter{enumi}{1}
\tightlist
\item
  我们置 \(\mu = \mu_d + \mu_a\)，其中 \(\mu_a\) 是原子测度，\(\mu_d\) 是弥散测度（INT，V，§5，第 10 节，命题 15）。设 S 为 \(\mu_a\) 的支集。我们有
\end{enumerate}

\[
\lim _ {\mathrm{H} \to + \infty} \frac {1}{2 \mathrm{H} + 1} \sum_ {| h | \leqslant \mathrm{H}} | \widehat {\mu} (h) | ^ {2} = \sum_ {x \in \mathrm{S}} | \mu (\{s \}) | ^ {2}
\]

（考察 \((\mu * \widetilde{\mu})(\{0\})\)）。

\begin{enumerate}
\def\labelenumi{\alph{enumi})}
\setcounter{enumi}{2}
\tightlist
\item
  设 \(\mu\) 为 \(\mathbf{T}\) 上的一个概率测度。下列条件等价：
\end{enumerate}

\begin{enumerate}
\def\labelenumi{(\roman{enumi})}
\item
  存在 \(x \in \mathbf{T}\) 使得 \(\mu = \varepsilon_x\)；
\item
  我们有 \(\lim_{|h|\to+\infty}|\widehat{\mu}(h)|^{2}=1\)；
\item
  我们有
\end{enumerate}

\[
\lim _ {\mathrm{H} \rightarrow + \infty} \frac {1}{2 \mathrm{H} + 1} \sum_ {| h | \leqslant \mathrm{H}} | \widehat {\mu} (h) | ^ {2} = 1.
\]

\begin{enumerate}
\def\labelenumi{\arabic{enumi})}
\setcounter{enumi}{64}
\tightlist
\item
  设 G 是紧的，并配备其规范化 Haar 测度 \(\mu\)。设 \(f: G \to G\) 为一个连续且满射的同态。
\end{enumerate}

\begin{enumerate}
\def\labelenumi{\alph{enumi})}
\item
  我们有 \(f(\mu) = \mu\)。
\item
  映射 \(f\) 是 \(\mu\)-遍历的（参看 I，第 190 页，练习 16）当且仅当不存在 G 的特征标 \(\chi \neq e_{\widehat{\mathrm{G}}}\) 与整数 \(n \geqslant 1\) 使得 \(\chi \circ f^n = \chi\)。
\item
  设 \(d \geqslant 1\) 为整数，a 为一个 \((d, d)\) 型的、行列式为 1 的整系数矩阵。映射 \(f_{a} \colon T^{d} \to T^{d}\)（由 \(f_{a}(x) = ax\) 对所有 \(x \in T^{d}\) 定义）关于 \(T^{d}\) 上的规范化 Haar 测度是遍历的，当且仅当矩阵 a 在 \(\mathcal{L}(\mathbf{C}^{n})\) 中的谱不含单位根。
\end{enumerate}

¶ 66) 我们用 \(e\colon \mathbf{T}\to \mathbf{C}^{*}\) 表示由特征标 \(x\mapsto e^{2i\pi x}\)（\(\mathbf{R}\) 的特征标）过渡到商而定义的特征标。

设 \(d \geqslant 2\) 为整数。设 \(\alpha \in R\)，并设 \(g \in R[X]\) 为一个次数 \textless{} d 的多项式。设 a 与 q 为互素的整数，其中 \(q \neq 0\)，使得

\[
\left| \alpha - \frac {a}{q} \right| \leqslant \frac {1}{q ^ {2}}.
\]

\begin{enumerate}
\def\labelenumi{\alph{enumi})}
\tightlist
\item
  设 \(\alpha \notin \mathbf{Z}\)。设 N 与 M 为满足 \(\mathrm{M} \geqslant 0\) 的整数。我们有
\end{enumerate}

\[
\Big | \sum_ {n = N} ^ {N + M} e ^ {2 i \pi \alpha n} \Big | \leqslant \frac {1}{| \sin (\pi \alpha) |}.
\]

\begin{enumerate}
\def\labelenumi{\alph{enumi})}
\setcounter{enumi}{1}
\tightlist
\item
  我们记 \(D = 2^{d-1}\)。对每个 \(\varepsilon > 0\)，存在实数 \(C(\varepsilon, d) \geqslant 0\)，使得对每个整数 \(P \geqslant 1\)，有
\end{enumerate}

\[
\left| \sum_ {n = 1} ^ {\mathrm{P}} e ^ {2 i \pi (\alpha n ^ {d} + g (n))} \right| \leqslant \mathrm{C} (\varepsilon , d) \mathrm{P} ^ {1 + \varepsilon} \left(\mathrm{P} ^ {- 1 / \mathrm{D}} + q ^ {- 1 / \mathrm{D}} + \left(\frac {\mathrm{P} ^ {d}}{q}\right) ^ {- 1 / \mathrm{D}}\right).
\]

（对 d 用归纳法，把不等式左端的平方化为次数 \(\leqslant d - 1\) 的一族多项式；当 d = 2 时，应用 a)。可以证明并利用估计式

\[
d (n) = \mathrm{O} (n ^ {\varepsilon})
\]

（当 \(n \to +\infty\)），其中 \(n \geqslant 1\)，\(d(n)\) 是 n 在 N 中的因子的个数，且此式对每个实数 \(\varepsilon > 0\) 都成立。

\begin{enumerate}
\def\labelenumi{\alph{enumi})}
\setcounter{enumi}{2}
\tightlist
\item
  由此给出练习 35 的 c) 中断言的一个新证明。
\end{enumerate}

¶ 67) 设 \(k \geqslant 2\) 与 \(\ell > k\) 为自然数。对每个整数 \(N \geqslant 1\)，我们用 \(r_{k,\ell}(N)\) 表示满足下式的族 \((n_1, \ldots, n_\ell) \in \mathbf{N}^\ell\) 的个数

\[
n _ {1} ^ {k} + \dots + n _ {\ell} ^ {k} = \mathrm{N}.
\]

本练习的目标是证明：若 \(\ell\) 相对于 \(k\) 足够大，则 \(r_{k,\ell}(\mathrm{N}) \geqslant 1\)（“Waring 问题”）对每个 \(\mathrm{N} \geqslant 1\) 成立。

我们用 \(e\colon\mathbf{T}\to\mathbf{C}^{*}\) 表示由特征标 \(x\mapsto e^{2i\pi x}\)（\(\mathbf{R}\) 的特征标）过渡到商而定义的特征标。

设 P ≥ 1 为整数。我们用 S \(_{P}\) 表示 T 上满足下式的函数

\[
\mathrm{S} _ {\mathrm{P}} (x) = \sum_ {n = 1} ^ {\mathrm{P}} e (x n ^ {k}).
\]

我们用 \(dx\) 表示 \(\mathbf{T}\) 上的规范化 Haar 测度。

\begin{enumerate}
\def\labelenumi{\alph{enumi})}
\tightlist
\item
  若 \(\mathrm{P} > \mathrm{N}^{1 / k}\)，则我们有
\end{enumerate}

\[
r _ {k, \ell} (\mathrm{N}) = \int_ {\mathbf {T}} \mathrm{S} _ {\mathrm{P}} (x) ^ {\ell} e (- \mathrm{N} x) d x.
\]

\begin{enumerate}
\def\labelenumi{\alph{enumi})}
\setcounter{enumi}{1}
\tightlist
\item
  设 \(\delta\) 为满足 \(0 < \delta < 1/3\) 的实数，并设 \(\mathrm{P} = [\mathrm{N}^{1/k}]\)。我们假设 N 充分大，使得 \(\mathrm{P} \geqslant 1\)。设 \(a\) 与 \(q\) 为满足 \(1 \leqslant q \leqslant \mathrm{P}^{\delta}\) 的互素整数。我们用 \(\mathrm{M}_{a,q}\) 表示开区间在 \(\mathbf{T}\) 中的像
\end{enumerate}

\[
\left. \right] \frac {a}{q} - \frac {1}{\mathrm{P} ^ {k - \delta}}, \frac {a}{q} + \frac {1}{\mathrm{P} ^ {k - \delta}} \Big [.
\]

对 \(x\in \mathrm{M}_{a,q}\)，有

\[
\sum_ {n = 1} ^ {\mathrm{P}} e ^ {2 i \pi x n ^ {k}} = \frac {1}{q} \mathrm{S} (a; q) \mathrm{I} \Bigl (x - \frac {a}{q} \Bigr) + \mathrm{O} (\mathrm{P} ^ {2 \delta})
\]

（当 N → +∞），其中

\[
\mathrm{S} (a; q) = \sum_ {x = 1} ^ {q} e \left(\frac {a x ^ {k}}{q}\right), \quad \mathrm{I} (y) = \int_ {0} ^ {\mathrm{P}} e \left(y t ^ {k}\right) d t.
\]

\begin{enumerate}
\def\labelenumi{\alph{enumi})}
\setcounter{enumi}{2}
\tightlist
\item
  我们用 M 表示诸集合 \(M_{a,q}\) 的并，其中 q 取遍满足 \(1 \leqslant q \leqslant P^{\delta}\) 的整数，a 取遍与 q 互素且满足 \(1 \leqslant a < q\) 的整数。若 \(\delta < 1/5\)，则存在实数 \(\delta' > 0\) 使得
\end{enumerate}

\[
\int_ {\mathrm{M}} \mathrm{S} _ {\mathrm{P}} (x) ^ {\ell} e (- \mathrm{N} x) d x = \mathrm{P} ^ {\ell - k} \mathfrak {S} (\mathrm{P} ^ {\delta}, \mathrm{N}) \mathrm{J} (\mathrm{P} ^ {\delta}) + \mathrm{O} (\mathrm{P} ^ {\ell - k - \delta^ {\prime}})
\]

（当 N → +∞），其中我们对 Q ≥ 1 置

\[
\mathfrak{S}(\mathrm{Q},\mathrm{N}) = \sum_{1\leqslant q\leqslant \mathrm{Q}}\sum_{\substack{a = 1\\ (a,q) = 1}}^{q}\Big(\frac{1}{q}\mathrm{S}(a;q)\Big)^{\ell}e\Big(-\frac{aN}{q}\Big)
\]

\[
\mathrm{J(Q)} = \int_ {- \mathrm{Q}} ^ {\mathrm{Q}} \Bigl (\int_ {0} ^ {1} e (y t ^ {k}) d t \Bigr) ^ {\ell} e (- y) d y.
\]

\begin{enumerate}
\def\labelenumi{\alph{enumi})}
\setcounter{enumi}{3}
\tightlist
\item
  我们有
\end{enumerate}

\[
\mathrm{J} (\mathrm{Q}) = \frac {\Gamma (1 + 1 / k) ^ {\ell}}{\Gamma (\ell / k)} + \mathrm{O} (\mathrm{Q} ^ {- (\ell / k - 1)})
\]

（当 Q → +∞）。（首先证明

\[
\mathrm{J} (\mathrm{Q}) = \frac {1}{k ^ {\ell}} \int_ {\mathrm{X} _ {\ell}} (y _ {1} \dots y _ {\ell - 1} (1 - y _ {1} - \dots - y _ {\ell - 1})) ^ {- 1 + 1 / k} d y + \mathrm{O} (\mathrm{Q} ^ {- (\ell / k - 1)})
\]

其中

\[
\mathrm{X} _ {\ell} = \left\{\left(y _ {1}, \dots , y _ {\ell - 1}\right) \in \left(\mathbf {R} _ {+} ^ {*}\right) ^ {\ell - 1} \mid y _ {1} + \dots + y _ {\ell - 1} <   1 \right\},
\]

然后计算这个积分；参看 FVR VII，第 8 页，命题 4。）

\begin{enumerate}
\def\labelenumi{\alph{enumi})}
\setcounter{enumi}{4}
\tightlist
\item
  设 \(\ell >\delta^{-1}k2^{k}\)。存在 \(\delta^{\prime} > 0\) 使得
\end{enumerate}

\[
\int_ {\mathbf {T - M}} | \mathrm{S} _ {\mathrm{P}} (x) | ^ {\ell} e (- \mathrm{N} x) d x = \mathrm{O} (\mathrm{P} ^ {\ell - k - \delta^ {\prime}})
\]

（当 \(N \rightarrow +\infty\)）。（若 N 充分大且 \(x \in R\) 的像位于 T-M 中，则存在互素的整数 a 与 q，使得

\[
\mathrm{P} ^ {\delta} \leqslant q \leqslant \mathrm{P} ^ {k - \delta}, \quad \left| x - \frac {a}{q} \right| \leqslant \frac {1}{q ^ {2}}.
\]

然后利用练习 66。）

\begin{enumerate}
\def\labelenumi{\alph{enumi})}
\setcounter{enumi}{5}
\tightlist
\item
  对每个素数 p 以及每个整数 \(m \geqslant 1\)，置
\end{enumerate}

\[
\mathrm{A}(p^{m}) = \sum_{\substack{a = 1\\ (a,q) = 1}}^{p^{m}}\Big(\frac{1}{p^{m}}\mathrm{S}(a;p^{m})\Big)^{\ell}e\Big(-\frac{a\mathrm{N}}{p^{m}}\Big)
\]

以及

\[
\mathrm{N} (p ^ {m}) = \operatorname{Card} \left(\left\{\left(x _ {1}, \dots , x _ {\ell}\right) \in (\mathbf {Z} / p ^ {m} \mathbf {Z}) ^ {\ell} \mid x _ {1} ^ {k} + \dots + x _ {\ell} ^ {k} = \mathrm{N mod.} p ^ {m} \right\}\right).
\]

当 \(\ell > 2^k\) 时，有

\[
1 + \sum_ {m = 1} ^ {+ \infty} \mathrm{A} (p ^ {m}) = \lim _ {m \to + \infty} \frac {\mathrm{N} (p ^ {m})}{p ^ {m (\ell - 1)}}.
\]

\begin{enumerate}
\def\labelenumi{\alph{enumi})}
\setcounter{enumi}{6}
\tightlist
\item
  当 \(\ell > 2^k\) 时，有 \(\mathfrak{S}(\mathrm{Q}, \mathrm{N}) = \mathfrak{S}(\mathrm{N}) + o(1)\)（当 \(\mathrm{Q} \to +\infty\)），其中
\end{enumerate}

\[
\mathfrak {S} (\mathrm{N}) = \prod_ {p} \Bigl (1 + \sum_ {m \geqslant 1} \mathrm{A} (p ^ {m}) \Bigr),
\]

该乘积取遍所有素数，且是绝对收敛的。

\begin{enumerate}
\def\labelenumi{\alph{enumi})}
\setcounter{enumi}{7}
\tightlist
\item
  设 p 为素数。设 \(\nu \geqslant 0\) 为使得 \(k = p^{\nu}r\)（其中 p 不整除 r）成立的整数。我们用 \(\gamma\) 表示整数
\end{enumerate}

\[
\gamma = \left\{ \begin{array}{l l} \nu + 1 & \text { if } p \geqslant 3 \\ \nu + 2 & \text { otherwise. } \end{array} \right.
\]

若存在整数 \((n_1, \ldots, n_\ell)\) 使得 \(p\) 不整除所有的 \(n_i\)，并且使得

\[
x _ {1} ^ {k} + \dots + x _ {\ell} ^ {k} \equiv 0 \mathrm{mod.} p ^ {\gamma},
\]

则

\[
\lim _ {m \to + \infty} \frac {\mathrm{N} (p ^ {m})}{p ^ {m (\ell - 1)}} > 0.
\]

\begin{enumerate}
\def\labelenumi{\roman{enumi})}
\tightlist
\item
  存在 \(\ell > k\) 使得：若 N 充分大，则 \(\mathfrak{S}(N) > 0\)。（利用上一小问以及练习 33 的 h)。）
\end{enumerate}

\begin{enumerate}
\def\labelenumi{\alph{enumi})}
\setcounter{enumi}{9}
\tightlist
\item
  存在整数 \(\ell > k\) 使得 \(r_{k,\ell}(\mathrm{N}) \geqslant 1\) 对每个整数 \(\mathrm{N} \geqslant 1\) 成立。
\end{enumerate}

\begin{enumerate}
\def\labelenumi{\arabic{enumi})}
\setcounter{enumi}{67}
\tightlist
\item
  我们把 G 的群运算写成加法。设 \(\Gamma\) 为 G 的一个离散子群，使得 \(G / \Gamma\) 是紧的。设 \(\Lambda = \Gamma^{\perp}\) 为 \(\Gamma\) 在 \(\widehat{G}\) 中的正交补；由于 \(\widehat{G / \Gamma}\) 与 \(\Lambda\) 等同（后者是 \(\widehat{G}\) 的离散子群），并且 \(\widehat{G} / \Lambda\) 是紧的（因为它的对偶与 \(\Gamma\) 等同）。
\end{enumerate}

我们假设 G 上的 Haar 测度如此选取，使得对偶测度 \(d\chi\) 在 \(\widehat{G}\) 上除以 \(\Lambda\) 上的计数测度所得之商，是 \(\widehat{G}/\Lambda\) 上的规范化 Haar 测度（INT，VII，§ 2，第 2 节，定义 1）。

\begin{enumerate}
\def\labelenumi{\alph{enumi})}
\item
  存在 \(\Omega\)（\(\widehat{\mathbf{G}}\) 的一个可测子集），它与 \(\mathbf{G}\) 模 \(\Lambda\) 的每个类恰好交于一点；\(\Omega\) 的测度等于 1。
\item
  对 \(x\in \mathbf{G}\)，置
\end{enumerate}

\[
\varphi (x) = \int_ {\Omega} \chi (x) d \chi .
\]

函数 \(\varphi\) 是连续的；它满足 \(\varphi(0) = 1\) 以及 \(\varphi(\gamma) = 0\)（对所有 \(\gamma \in \Gamma - \{0\}\)）。

\begin{enumerate}
\def\labelenumi{\alph{enumi})}
\setcounter{enumi}{2}
\tightlist
\item
  证明 \(\varphi\) 属于 \(\mathrm{L}^2 (\mathrm{G})\) 且 \(\| \varphi \| _2 = 1\)。对每个 \(\gamma \in \Gamma -\{e\}\)，有
\end{enumerate}

\[
\int_ {\mathrm{G}} \varphi (x) \overline {{\varphi (\gamma - x)}} d x = 0.
\]

\begin{enumerate}
\def\labelenumi{\alph{enumi})}
\setcounter{enumi}{3}
\tightlist
\item
  设 \(f \in \mathrm{L}^{2}(\mathrm{G})\) 满足：\(\mathcal{F}(f)\) 在 \(\Omega\) 之外几乎处处为零。证明
\end{enumerate}

\[
f = \sum_ {\gamma \in \Gamma} f (\gamma) \boldsymbol {\delta} _ {\gamma} (\varphi)
\]

（在 \(\mathrm{L}^2 (\mathrm{G})\) 中），其中 \(\delta_{\gamma}(\varphi)(x) = \varphi (x - \gamma)\)。此外，我们有

\[
\int_ {\mathrm{G}} | f (x) | ^ {2} d x = \sum_ {\gamma \in \Gamma} | f (\gamma) | ^ {2}.
\]

\begin{enumerate}
\def\labelenumi{\alph{enumi})}
\setcounter{enumi}{4}
\tightlist
\item
  存在 G 上唯一的连续函数，它与 f 几乎处处相等。若 f 是连续的，则
\end{enumerate}

\[
f = \sum_ {\gamma \in \Gamma} f (\gamma) \boldsymbol {\delta} _ {\gamma} (\varphi)
\]

（在 \(\mathcal{C}(\mathrm{G})\) 中）。

\begin{enumerate}
\def\labelenumi{\alph{enumi})}
\setcounter{enumi}{5}
\tightlist
\item
  特殊化到 G = R 且 \(\Lambda = hZ\)（其中 h \textgreater{} 0）的情形，取 \(\Omega = [-h, h[\)。
\end{enumerate}

\exercisesection{分类}

\begin{enumerate}
\def\labelenumi{\arabic{enumi})}
\tightlist
\item
  \begin{enumerate}
  \def\labelenumii{\alph{enumii})}
  \tightlist
  \item
    为使 G 具有可数基，必要且充分的是 \(\widehat{G}\) 具有可数基。（若 G 有可数基，则 \(L^{1}(G)\) 有可数基，从而 \(\mathsf{X}(\mathsf{L}^{1}(\mathsf{G}))\) 有可数基。）
  \end{enumerate}
\end{enumerate}

\begin{enumerate}
\def\labelenumi{\alph{enumi})}
\setcounter{enumi}{1}
\item
  为使 G 在无穷远处可数，必要且充分的是 \(\widehat{\mathrm{G}}\) 容有一个具有可数基的开子群，或者等价地，\(\widehat{\mathrm{G}}\) 是可度量化的。（利用 II，第 248 页，命题 3。）
\item
  设 G 是紧的。为使 G 可度量化，必要且充分的是 \(\widehat{G}\) 是可数的，或者等价地，G 同构于 \(T^{N}\) 的一个闭子群。（一个可数的离散交换群是 \(\mathbf{Z}^{(N)}\) 的一个商。）
\end{enumerate}

\begin{enumerate}
\def\labelenumi{\arabic{enumi})}
\setcounter{enumi}{1}
\tightlist
\item
  \begin{enumerate}
  \def\labelenumii{\alph{enumii})}
  \tightlist
  \item
    每个无限交换群 \(\Gamma\) 都容许一个无限的可数商群。（利用 A，II，第 185 页，练习 14，把一个无限可数子群 \(\Delta\)（\(\Gamma\) 的）嵌入到一个可数的可除群 \(\Delta'\) 中；然后在 \(\Gamma\) 上延拓 \(\Delta\) 到 \(\Delta'\) 的恒等同态。）
  \end{enumerate}
\end{enumerate}

\begin{enumerate}
\def\labelenumi{\alph{enumi})}
\setcounter{enumi}{1}
\tightlist
\item
  每个无限的紧交换群都包含一个无限的可度量化闭子群。（利用 a)。）
\end{enumerate}

\begin{enumerate}
\def\labelenumi{\arabic{enumi})}
\setcounter{enumi}{2}
\tightlist
\item
  \begin{enumerate}
  \def\labelenumii{\alph{enumii})}
  \tightlist
  \item
    设 K 为紧群，\(\mathrm{K}_{(n)}\) 为 K 中阶为 n 的元素所成的集合。若对每个 n 都有 \(K \neq E_{n}\)，则 K 中无限阶元素所成的集合在 K 中稠密。（若 \(\mathrm{K}_{(n)}\) 有非空内部，则有 \(\mathrm{K}_{(m)} = \mathrm{K}\)（对某个整数 \(m \geqslant n\)）。）
  \end{enumerate}
\end{enumerate}

\begin{enumerate}
\def\labelenumi{\alph{enumi})}
\setcounter{enumi}{1}
\item
  设 \(\Gamma\) 为一个离散交换群，其元素的阶不是有界的。则 \(\Gamma\) 容许一个具有同样性质的可数商群。（按练习 2 的 a) 同样论证。）
\item
  若 e 在 G 中的每个邻域都包含一个无限阶元素，则 G 具有一个带同样性质的可度量化闭子群。（化归到 G 是紧的情形，并利用 a) 与 b)。）反之，若 G 不是离散的，则存在整数 q \textgreater{} 1，使得 G 包含一个同构于 \((\mathbf{Z}/q\mathbf{Z})^{\mathbf{N}}\) 的闭子群。（利用 A，VII，第 54 页，练习 4，c)。）
\end{enumerate}

\begin{enumerate}
\def\labelenumi{\arabic{enumi})}
\setcounter{enumi}{3}
\item
  若 G 是紧的且 \(\widehat{G}\) 有一个无限阶元素，则存在一个从 R 到 G 的非平凡连续同态。（由于 R 是可除的，存在一个从 \(\widehat{G}\) 到 R 的非平凡同态。）
\item
  设 p 为素数。群 \(Q_{p}\) 不是紧群与离散群的乘积。（\(Q_{p}\) 的每个紧子群都具有形式 \(p^{n}Z_{p}\)（对某个整数 \(n \in Z\)）。）
\item
  若 G 是挠群，则 G 与 \(\widehat{G}\) 都是全不连通的。（由 II，第 250 页，推论 2，\(\widehat{G}\) 是全不连通的；为证 G 全不连通，利用 II，第 248 页，命题 3 化归到 G 为紧的情形，然后证明 \(\widehat{G}\) 是挠群。）
\item
  我们称元素 \(x\)（属于 \(G\)）为无限高的，如果对每个 \(n \in \mathbf{Z}\)，存在 \(y \in G\) 使得 \(y^n = x\)。证明：若 \(G\) 是紧的，则无限高元素所成的集合是 \(G\) 的中性分支。
\end{enumerate}

¶ 8) 群 G 是局部连通的，当且仅当 G 同构于一个形如 \(\mathbf{R}^{n}\times\mathrm{E}\times\widehat{\mathrm{D}}\) 的群，其中 E 与 D 是离散的，且 D 的每个有限秩子群都是自由的。

\begin{enumerate}
\def\labelenumi{\arabic{enumi})}
\setcounter{enumi}{8}
\tightlist
\item
  \begin{enumerate}
  \def\labelenumii{\alph{enumii})}
  \tightlist
  \item
    设 \(\boldsymbol{a}=(a_{n})_{n\geqslant0}\) 是一个整数序列，均 \textgreater{} 1。我们给 Z 赋予如下的拓扑环结构：其中子集 \(Za_{0}a_{1}\cdots a_{n}\)（\(n\geqslant0\)）构成 0 的一个基本邻域系。我们用 \(\Delta_{a}\) 表示这个环的完备化。它是紧的、可度量的、全不连通的。若 p 是素数且对所有 i 有 \(a_{i}=p\)，则我们有 \(\Delta_{a}=Z_{p}\)。
  \end{enumerate}
\end{enumerate}

\begin{enumerate}
\def\labelenumi{\alph{enumi})}
\setcounter{enumi}{1}
\item
  我们用 \(\mathbf{Z}(\boldsymbol{a}^{\infty})\) 表示形如 \(\exp(2i\pi l/(a_{0}a_{1}\cdots a_{r}))\)（其中 \(l \in Z\) 且 \(r \in N\)）的复数所成的乘法群，并赋予离散拓扑。则 \(\chi \mapsto \chi(1)\) 是从 \(\widehat{\Delta}_{\boldsymbol{a}}\) 到 \(\mathbf{Z}(\boldsymbol{a}^{\infty})\) 上的同构。
\item
  设 B 为 \(\mathbf{R} \times \Delta_{\boldsymbol{a}}\) 的由 \((n, n)\)（\(n \in \mathbf{Z}\)）组成的子群；它是离散的。我们置 \(\Sigma_{\boldsymbol{a}} = (\mathbf{R} \times \Delta_{\boldsymbol{a}})/\mathrm{B}\)。它是一个紧的、可度量的连通群。（注意 \(\mathbf{R}\) 在 \(\Sigma_{\boldsymbol{a}}\) 中的典范像是稠密的，并利用 II，第 244 页，引理 1。）
\end{enumerate}

\(d)\) 设 \(\Gamma_{\boldsymbol{a}}\) 为离散群 \(\mathbf{Q}\) 的由形如 \(l/(a_{0}a_{1}\cdots a_{r})\)（其中 \(l\in\mathbf{Z}\) 且 \(r\in\mathbf{N}\)）的有理数组成的加法子群。若 \(\chi\in\widehat{\Sigma}_{\boldsymbol{a}}\)，则 \(\chi\) 定义了 \(\mathbf{R}\times\Delta_{\boldsymbol{a}}\) 的一个酉特征标，从而定义了 \(\mathbf{R}\) 的一个形如 \(x\mapsto\exp(2i\pi\alpha_{\chi}x)\) 的酉特征标（其中 \(\alpha_{\chi}\in\mathbf{R}\)）。则 \(\chi\mapsto\alpha_{\chi}\) 是从 \(\widehat{\Sigma}_{\boldsymbol{a}}\) 到 \(\Gamma_{\boldsymbol{a}}\) 上的同构。

\begin{enumerate}
\def\labelenumi{\alph{enumi})}
\setcounter{enumi}{4}
\item
  证明 \(\widehat{\mathbf{Q}}\) 典范地同构于 \(\Sigma_{\pmb{a}}\)，其中序列 \(\pmb{a}\) 由 \(a_{n} = n + 2\)（对所有 \(n \geqslant 0\)）定义。
\item
  设 \(R_{d}\) 为赋予离散拓扑的群 R。证明 \(\widehat{R}_{d}\) 同构于 \((\widehat{\Sigma}_{\boldsymbol{a}})^{\mathfrak{c}}\)，其中 \(a_{n}=n+2\)，c 是连续统的基数。（考察 R 作为 Q-向量空间的一组基。）
\end{enumerate}

¶ 10) 称一个拓扑群为单生的，如果存在群中的一个元素，它的幂在群中稠密；称它为螺线管群，如果存在从 R 到该群中的连续同态，其像是稠密的。

\begin{enumerate}
\def\labelenumi{\alph{enumi})}
\item
  沿用练习 9 的记号，群 \(\Delta_{\pmb{a}}\) 是单生的。
\item
  每个紧的单生全不连通群都拓扑同构于形如 \(\Delta_{a}\) 的群。
\item
  设 G 是紧的。为使 G 是单生的，必要且充分的是 \(\widehat{G}\) 同构于赋予离散拓扑的 T 的一个子群。
\item
  设 G 是紧的。下列条件等价：
\end{enumerate}

\begin{enumerate}
\def\labelenumi{(\roman{enumi})}
\item
  G 是螺线管群；
\item
  \(\widehat{\mathbf{G}}\) 同构于赋予离散拓扑的 \(\mathbf{R}\) 的一个子群；
\item
  \(\widehat{\mathbf{G}}\) 是无挠的，且 \(\operatorname{Card}(\widehat{\mathbf{G}}) \leqslant \mathfrak{c}\)（连续统的基数）；
\item
  G 是 \((\Sigma_{\boldsymbol{a}})^{\mathfrak{c}}\) 的一个商，其中 \(\boldsymbol{a}=(n+2)_{n\geqslant0}\)。
\end{enumerate}

\begin{enumerate}
\def\labelenumi{\alph{enumi})}
\setcounter{enumi}{4}
\tightlist
\item
  若 G 是紧的螺线管群，则 G 是单生的。
\end{enumerate}

\begin{enumerate}
\def\labelenumi{\arabic{enumi})}
\setcounter{enumi}{10}
\tightlist
\item
  我们用 L(G) 表示从 G 到 R 中的连续表示所成的集合。G 的单参数子群是指从 R 到 G 中的连续同态的像。
\end{enumerate}

\begin{enumerate}
\def\labelenumi{\alph{enumi})}
\tightlist
\item
  下列条件等价：
\end{enumerate}

\begin{enumerate}
\def\labelenumi{(\roman{enumi})}
\item
  G 是诸单参数子群的并；
\item
  从 Z 到 G 的每个同态都可延拓为从 R 到 G 的连续同态；
\item
  从 \(\widehat{\mathbf{G}}\) 到 \(\mathbf{R} / \mathbf{Z}\) 的每个连续同态都由 \(\mathrm{L}(\widehat{\mathrm{G}})\) 中的一个元素过渡到商而得到；
\item
  \(\widehat{\mathbf{G}}\) 的每个酉特征标都形如 \(e^{i\lambda}\)（其中 \(\lambda \in \mathrm{L}(\widehat{\mathbf{G}})\)）。
\end{enumerate}

\begin{enumerate}
\def\labelenumi{\alph{enumi})}
\setcounter{enumi}{1}
\item
  若 G 有可数基，则 a) 中的条件等价于：G 同构于 \(R^{n} \times T^{I}\)（其中 I 是一个可数集）。
\item
  G 的诸单参数子群的并是 G 的中性分支的一个稠密子群。（利用练习 4。）
\end{enumerate}

\(d)\) 设 \(t \mapsto \chi_t\) 是从 \([0,1]\) 到 \(\widehat{\mathbf{G}}\) 中的连续映射，使得 \(\chi_0\) 是平凡特征标。存在一个映射 \(t \mapsto \lambda_t\)（从 \([0,1]\) 到 \(\mathrm{L}(\mathrm{G})\)，满足 \(\lambda_0 = 0\)），使得对所有 \(t\) 有 \(\chi_t = \exp(2i\pi\lambda_t)\)，并且对每个 \(x \in \mathrm{G}\)，映射 \(t \mapsto \lambda_t(x)\) 是连续的。

\begin{enumerate}
\def\labelenumi{\alph{enumi})}
\setcounter{enumi}{4}
\tightlist
\item
  G 的诸单参数子群的并也是 G 的包含 e 的弧的并，G 的弧是指从 \([0,1]\) 到 G 中的连续映射的像。（利用 d)。
\end{enumerate}

\begin{enumerate}
\def\labelenumi{\arabic{enumi})}
\setcounter{enumi}{11}
\tightlist
\item
  我们沿用练习 11 中的记号 L(G)。
\end{enumerate}

\begin{enumerate}
\def\labelenumi{\alph{enumi})}
\item
  设 H 为 G 的闭子群。L(H) 的每个元素都可延拓为 L(G) 的元素。（先化归到 \(G = R^{n} \times D\)（D 离散）的情形，再化归到 L(H) 的给定元素在 \(R^{n}\) 上为平凡的情形，然后化归到 G 为离散的情形。最后利用 R 是可除的这一事实。）
\item
  从 H 到 \(R^{n} \times T^{I}\)（其中 I 是任意集合）的每个连续态射都可延拓为从 G 到 \(R^{n} \times T^{I}\) 的连续态射。（利用 a) 以及 II，第 226 页，定理 4。）
\item
  反之，设 A 为一个局部紧交换群。假设对每个局部紧交换群 G、对 G 的每个闭子群 H 以及从 H 到 A 的每个连续态射 \(\varphi\)，\(\varphi\) 都可延拓为从 G 到 A 的连续态射。则存在整数 n 与集合 I，使得 A 同构于 \(R^{n} \times T^{I}\)。（取 G = R，H = Z，可以看出 A 是连通的，从而形如 \(R^{n} \times \widehat{D}\)（D 离散）。证明 D 是一个投射 Z-模，从而是自由的。）
\end{enumerate}

\begin{enumerate}
\def\labelenumi{\arabic{enumi})}
\setcounter{enumi}{12}
\tightlist
\item
  设 C 为 G 的紧子集，\(\varepsilon > 0\)。
\end{enumerate}

\begin{enumerate}
\def\labelenumi{\alph{enumi})}
\item
  设 \(\alpha\) 为 G 的 Haar 测度。存在 G 的紧子集 D，使得 \(\alpha(\mathrm{CD}) \leqslant (1 + \varepsilon)\alpha(\mathrm{D})\)。（\(\mathbf{G} = \mathbf{R}^n\) 的情形是直接的，我们化归到 G 容有一个紧开子群 H 的情形，并且可以假设 C 关于 H 是饱和的；于是化归到 G 为离散的情形；接着可以假设 G 是有限生成的，最后 \(\mathbf{G} = \mathbf{Z}^p\)。）
\item
  存在 \(k \in \mathrm{L}^{1}(\widehat{\mathrm{G}})\) 使得 \(\mathcal{F}(k) \in \mathcal{K}(\mathrm{G})\) 且 \(\mathcal{F}(k)(x) = 1\)（对 \(x \in C\)），并有 \(\|k\|_{1} \leqslant 1 + \varepsilon\)。（取形如 \(\lambda fg\) 的 k，其中 \(\lambda\) 是常数，\(f, g \in \mathrm{L}^{2}(\widehat{\mathrm{G}})\) 的 Fourier 变换是适当集合的特征函数，利用 a)。）
\end{enumerate}

¶ 14) a) 假设 e 在 G 中的每个邻域都包含一个无限阶元素。存在 G 的一个紧的、可度量的、完满的全不连通子集 P，它是一个 Kronecker 集（II，第 265 页，练习 15），并且存在一个集中在 P 上的非零弥散测度。（利用练习 3 的 c) 化归到 G 可度量化的情形。然后仿照 TG，IX，第 64 页，引理 3 的构造，并同时利用 II，第 265 页，练习 15 的 e)。）

\begin{enumerate}
\def\labelenumi{\alph{enumi})}
\setcounter{enumi}{1}
\item
  若 \(\mathrm{G} = (\mathbf{Z}/q\mathbf{Z})^{\mathbf{N}}\)，则存在 G 的一个紧的完满全不连通子集 P，它是 \(K_{q}\) 型的，并且存在一个集中在 P 上的非零弥散测度。
\item
  设 G 不是离散的。存在元素 \(\tau\)（属于 \(\mathcal{M}^{1}(G)\)）使得 \(\|\mathcal{F}(\tau)\|_{\infty}<\varrho(\tau)\)，其中 \(\varrho(\tau)\) 是 \(\mathcal{M}^{1}(G)\) 中的谱半径。存在不可逆元素 \(\tau'\)（属于 \(\mathcal{M}^{1}(G)\)）使得 \(\mathcal{F}(\tau')\geqslant1\) 在 \(\widehat{G}\) 上处处成立。存在 \(\mathcal{M}^{1}(G)\) 的一个非 Hermite 特征标。（利用 a)、b)、练习 3 的 c)，以及 II，第 265 页，练习 15 的 f)。）
\end{enumerate}

\(d)\) 由 \(c\) ) 推出：\(\widehat{\mathrm{G}}\)（它在 \(\mathsf{X}(\mathcal{M}^1 (\mathrm{G}))\) 中是开的）在 \(\mathsf{X}(\mathcal{M}^1 (\mathrm{G}))\) 中不是稠密的。

\exercisesection{不变子空间}

\begin{enumerate}
\def\labelenumi{\arabic{enumi})}
\tightlist
\item
  设 \((a_{n})_{n\in \mathbf{N}}\) 为一列复数，满足：存在实数 \(A\geqslant 0\)，使得 \(n|a_n|\leqslant A\)（对所有 \(n\in \mathbf{N}\)）。对 \(z\in \mathbf{C}\)（模 \(<  1\)），置 \(\varphi (z) = \sum_{n\geqslant 0}a_nz^n\)。
\end{enumerate}

我们假设存在 \(\ell \in \mathbf{C}\) 使得

\[
\lim_{\substack{x\in \mathbf{R}_{+}\\ x <   1}}\varphi (x) = \ell .
\]

对 \(t\in \mathbf{R}_+^*\)，置

\[
g (t) = \sum_ {0 \leqslant n \leqslant t ^ {- 1}} a _ {n}, \quad \mathrm{F} (t) = \varphi (e ^ {- t}).
\]

我们用 \(dx\) 表示 \(\mathbf{R}\) 上的 Lebesgue 测度，以及它在 \(\mathbf{R}_{+}^{*}\) 上的限制。

\begin{enumerate}
\def\labelenumi{\alph{enumi})}
\tightlist
\item
  设 \(s \leqslant t\) 为 \(\mathbf{R}_{+}^{*}\) 中的元素。我们有
\end{enumerate}

\[
| g (s) - g (t) | \leqslant | a _ {[ t ^ {- 1} + 1 ]} | + A \log \left(\frac {t}{s}\right).
\]

\begin{enumerate}
\def\labelenumi{\alph{enumi})}
\setcounter{enumi}{1}
\item
  函数 g 当 \(t \rightarrow 0\) 时是缓变的。
\item
  对每个 \(t \in R_{+}^{*}\)，函数 \(x \mapsto g(t/x)xe^{-x}\) 在 \(R_{+}^{*}\) 上关于 Haar 测度 \(x^{-1}dx\) 是可积的。
\item
  对 \(0 < x < 1\)，有
\end{enumerate}

\[
\varphi (x) = \sum_ {n \geqslant 0} g \left(\frac {1}{n}\right) \left(x ^ {n} - x ^ {n + 1}\right).
\]

\begin{enumerate}
\def\labelenumi{\alph{enumi})}
\setcounter{enumi}{4}
\tightlist
\item
  设 \(f(x) = xe^{-x}\)（对 \(x \in R_{+}^{*}\)）。我们有 \(f \in \mathrm{L}^{1}(\mathbf{R}_{+}^{*}, dx/x)\)；它的积分等于 1，并且我们有 F = g * f。
\end{enumerate}

\(f)\) 函数 \(g\) 属于 \(\mathrm{L}^{\infty}(\mathbf{R}_{+}^{*})\)。（证明 \(\mathrm{F}-g\) 是有界的（利用 \(a\) ) 以及假设）。）

\(g)\) 函数 \(f\) 在群 \(\mathbf{R}_{+}^{*}\) 上的 Fourier 变换不取零值。（把 \(\mathcal{F}(f)\) 等同于函数 \(y\mapsto\Gamma(1+iy)\) 并利用 FVR, VII。）

\begin{enumerate}
\def\labelenumi{\alph{enumi})}
\setcounter{enumi}{7}
\tightlist
\item
  我们有
\end{enumerate}

\[
\sum_ {n = 0} ^ {+ \infty} a _ {n} = \ell .
\]

这个命题就是 Littlewood 的“Tauber 型”定理。关于序列 \((a_n)\) 的假设，即 \(a_n = \mathrm{O}(1/n)\)，还可以减弱：只需存在常数 \(A \geqslant 0\)，使得 \(na_n \geqslant -A\)（对每个整数 \(n \geqslant 1\)）（“Hardy–Littlewood Tauber 型定理”，例如参看 J. KOREVAAR, Tauberian theory, Grundlehren math. wiss. 329, Springer (2004), Th. 7.4, 第 16 页；读者应注意不要把这个定理与 FVR, I, 第 50 页, 练习 18 中的定理混淆）。

\begin{enumerate}
\def\labelenumi{\arabic{enumi})}
\setcounter{enumi}{1}
\item
  设 F 为 \(\widehat{\mathbf{G}}\) 的闭子集。设 K 为 \(\widehat{\mathbf{G}}\) 的紧子集，使得 \(\mathrm{F} \cap \mathrm{K} = \varnothing\)。证明存在 G 上可积的连续函数 \(g\)，使得 \(\mathcal{F}(g)\) 在 F 上等于 0，在 K 上等于 1。（设 \(f \in \mathrm{L}^{1}(\mathrm{G})\) 满足 \(\mathcal{F}(f)\) 在 F 上为零且在 K 上等于 1；考察 \(f * \varphi\)，其中 \(\varphi\) 是形如 \(h * h'\) 的函数的 Fourier 变换（\(h \in \mathcal{K}(\widehat{\mathrm{G}})\)，\(h' \in \mathcal{K}(\widehat{\mathrm{G}})\)，且 \(h * h' = 1\) 在 K 上）。）
\item
  设 \(g\) 是定义在 \(]0, +\infty[\) 上的复值函数，关于 Lebesgue 测度可积。我们假设 \(\int_0^{+\infty} g(t)dt = 1\)，且 \(\int_0^{+\infty} g(t)t^{ix}dt \neq 0\)（对所有 \(x \in \mathbf{R}\)）。设 \(f\) 是 \(]0, +\infty[\) 上的复值可测有界函数，使得 \(\frac{1}{x}\int_0^{+\infty} g(t/x)f(t)dt\) 趋于一个有限极限 \(\ell\)（当 \(x\) 趋于 \(+\infty\)）。
\end{enumerate}

\begin{enumerate}
\def\labelenumi{\alph{enumi})}
\tightlist
\item
  对定义在 \(]0, +\infty[\) 上的、关于 Lebesgue 测度可积且积分为 1 的任意复值函数 h，我们有
\end{enumerate}

\[
\lim _ {x \to + \infty} \frac {1}{x} \int_ {0} ^ {+ \infty} h \left(\frac {t}{x}\right) f (t) d t = \ell .
\]

特别地，

\[
\lim _ {x \rightarrow + \infty} \frac {1}{x} \int_ {0} ^ {x} f (t) d t = \ell .
\]

\((\text{设 } g_1(t) = tg(t), h_1(t) = th(t), \text{ 我们有 } g_1 \in \mathrm{L}^1(\mathbf{R}_+^*) ; \text{ Fourier 变换 } \mathcal{F}(g_1) \text{ 不取零值且 } \check{g}_1 * f \text{ 趋于 } \ell, \text{ 因此 } h_1 * f \text{ 趋于 } \ell.)\) b) 若 \(f(t)\) 在 \(\mathbf{R}_+^*\) 上是缓变的（当 \(t\) 趋于 \(+\infty\)），则 \(f\) 趋于 \(\ell\)（当 \(t\) 趋于 \(+\infty\)）。

\begin{enumerate}
\def\labelenumi{\arabic{enumi})}
\setcounter{enumi}{3}
\tightlist
\item
  设 G = R。对每个实数 \(\alpha > 1\)，我们用下式定义 \(f_{\alpha} \in L^{1}(\mathbf{R})\)
\end{enumerate}

\[
f _ {\alpha} (x) = \left\{ \begin{array}{l l} 2 & \text { if } 0 <   x <   1 \\ 1 & \text { if } 1 \leqslant x <   \alpha \\ 0 & \text { otherwise. } \end{array} \right.
\]

\begin{enumerate}
\def\labelenumi{\alph{enumi})}
\item
  函数 \(f * \varepsilon_{x}\)（\(x \in R\)）在 \(\mathrm{L}^{1}(\mathbf{R})\) 中构成完全集，当且仅当 \(\alpha\) 是无理数。
\item
  函数 \(f * \varepsilon_x\)（\(x \in \mathbf{R}\)）在 \(L^2(\mathbf{R})\) 中构成完全集。
\end{enumerate}

\begin{enumerate}
\def\labelenumi{\arabic{enumi})}
\setcounter{enumi}{4}
\tightlist
\item
  设 \(f \in \mathrm{L}^{\infty}([0, +\infty([)\text{ 且 } k \in \mathbf{R}_{+}^{*}. \text{ 若}\)
\end{enumerate}

\[
\lim _ {x \to + \infty} \frac {1}{x} \int_ {0} ^ {+ \infty} e ^ {- t / x} f (t) d t = \ell
\]

则

\[
\lim _ {x \rightarrow + \infty} \frac {k}{x ^ {k}} \int_ {0} ^ {x} (x - t) ^ {k - 1} f (t) d t = \ell
\]

（利用函数 \(g(t) = e^{-t}\)，以及 \(h(t) = k(1 - t)^{k-1}\)（对 \(0 \leqslant t \leqslant 1\)；\(h(t) = 0\) 对 \(t > 1\)）。）

\begin{enumerate}
\def\labelenumi{\arabic{enumi})}
\setcounter{enumi}{5}
\tightlist
\item
  设 \(f \colon \mathbf{R} \to \mathbf{R}\) 与 \(g: \mathbf{R} \to \mathbf{C}\) 是两个函数。我们假设 \(f \in \mathrm{L}^{\infty}(\mathbf{R})\) 且 \(g \in \mathrm{L}^{1}(\mathbf{R})\)，\(\int_{\mathbf{R}} g(x) dx = 1\)，并且 \(\mathcal{F}(g)\) 不取零值。我们假设 \(f\) 是缓减的，即：对满足 \(y > x\)（\(y - x\) 趋于 0 且 \(x \to +\infty\)）的 y、x，有
\end{enumerate}

\[
\liminf (f (y) - f (x)) \geqslant 0.
\]

若 \((g*f)(x)\) 趋于 \(\ell\)（当 x 趋于 \(+\infty\)），则 \(f(x)\) 趋于 \(\ell\)（当 x 趋于 \(+\infty\)）。

\begin{enumerate}
\def\labelenumi{\arabic{enumi})}
\setcounter{enumi}{6}
\tightlist
\item
  设 \(g: R \rightarrow C\) 为连续函数，满足：
\end{enumerate}

\[
\sum_ {n = - \infty} ^ {+ \infty} \sup _ {n \leqslant t \leqslant n + 1} | g (t) | <   + \infty .
\]

\begin{enumerate}
\def\labelenumi{\alph{enumi})}
\item
  我们有 \(g \in \mathrm{L}^{1}(\mathbf{R})\)。
\item
  设 g 的积分为 1 且 \(\mathcal{F}(g)\) 不取零值。设 \(\mu\) 为 R 上的复测度，使得 \(|\mu|([x,x+1])\) 当 x 取遍 R 时是有界的。若 \((g*\mu)(x)\) 趋于 \(\ell\)（当 x 趋于 \(+\infty\)），则对每个满足下述条件的连续函数 \(h:R\to C\)
\end{enumerate}

\[
\sum_ {n = - \infty} ^ {+ \infty} \sup _ {n \leqslant t \leqslant n + 1} | h (t) | <   + \infty ,
\]

（积分为 1），\((h * \mu)(x)\) 当 \(x\) 趋于 \(+\infty\) 时的极限等于 \(\ell\)。（证明 \(h * \mu\) 是缓变的，并且 \((g * (h * \mu))(x)\) 趋于 \(\ell\)（当 \(x\) 趋于 \(+\infty\)）。）

\begin{enumerate}
\def\labelenumi{\arabic{enumi})}
\setcounter{enumi}{7}
\tightlist
\item
  设 \(g\) 是满足练习 3 的假设的函数。此外，设 \(g \geqslant 0\) 且 \(\liminf_{x \to 0} g(x) > 0\)。设 \(f\) 是 \(]0, +\infty[\) 上的有下界的实值可测函数。若
\end{enumerate}

\[
\lim _ {x \to + \infty} \frac {1}{x} \int_ {0} ^ {+ \infty} g \left(\frac {t}{x}\right) f (t) d t = \ell
\]

则

\[
\sigma (x) = \frac {1}{x} \int_ {0} ^ {x} f (t) d t
\]

趋于 \(\ell\)（当 \(x\) 趋于 \(+\infty\)）。（证明 \(\sigma\) 对 \(x \geqslant 1\) 是有界的，再证明 \(\sigma\) 是缓减的，然后证明

\[
\frac {1}{x} \int_ {0} ^ {+ \infty} g \left(\frac {t}{x}\right) \sigma (t) d t
\]

趋于 \(\ell\)（当 \(x\) 趋于 \(+\infty\)）。）

\begin{enumerate}
\def\labelenumi{\arabic{enumi})}
\setcounter{enumi}{8}
\tightlist
\item
  设 \(\varphi\colon[0,+\infty[\to\mathbf{R}\) 是非负的递增函数，使得函数 \(t\mapsto e^{-\sigma t}\varphi(t)\) 当 \(\sigma>1\) 时关于 Lebesgue 测度可积。对每个 \(s\in\mathbf{C}\)（满足 \(\mathcal{R}(s)>1\)），置
\end{enumerate}

\[
f (s) = \int_ {0} ^ {+ \infty} e ^ {- s t} \varphi (t) d t
\]

（“\(\varphi\) 的 Laplace 变换”）。

假设存在 \(A \in C\) 具有下列性质：当 \(\sigma\) 通过 \textgreater{} 1 的值趋于 1 时，函数

\[
t \mapsto f (\sigma + i \tau) - \frac {\mathrm{A}}{\sigma + i \tau - 1} \qquad \qquad \tau \in \mathbf {R}
\]

在 R 的每个紧子集上一致收敛于一个函数 g。则

\[
\lim _ {t \to + \infty} \varphi (t) e ^ {- t} = A
\]

（“Ikehara Tauber 型定理”）。

（我们置 \(a(t) = e^{-t}\varphi(t)\)（对 \(t > 0\)），\(a(t) = 0\)（对 \(t \leqslant 0\)），\(A(t) = A\)（对 \(t > 0\)），\(A(t) = 0\)（对 \(t \leqslant 0\)）。设 \(\lambda > 0\)。我们置：

\[
k _ {\lambda} (t) = \lambda \sqrt {\frac {2}{\pi}} \Big (\frac {\sin 2 \lambda t}{2 \lambda t} \Big) ^ {2}
\]

\[
\mathrm{K} _ {\lambda} (t) = \left\{ \begin{array}{l l} 1 - \left| \frac {t}{2 \lambda} \right| & \text { if } | t | \leqslant 2 \lambda \\ 0 & \text { if } | t | > 2 \lambda . \end{array} \right.
\]

利用 II，第 262 页，练习 1，Lebesgue–Fubini 定理以及 Plancherel 定理，证明：

\[
\lim _ {\varepsilon \to 0} \frac {1}{\sqrt {2 \pi}} \int_ {- \infty} ^ {+ \infty} k _ {\lambda} (x - t) (a (t) - \mathrm{A} (t)) e ^ {- \varepsilon t} d t = \frac {1}{\sqrt {2 \pi}} \int_ {- 2 \lambda} ^ {2 \lambda} \mathrm{K} _ {\lambda} (y) e ^ {- i x y} g (y) d y
\]

并推出

\[
\lim _ {x \to + \infty} \frac {1}{\sqrt {2 \pi}} \int_ {- \infty} ^ {+ \infty} k _ {\lambda} (x - t) a (t) d t = \mathrm{A}.
\]

然后证明 \(a\) 是有界的且缓减的。最后可以应用练习 6。）

\begin{enumerate}
\def\labelenumi{\arabic{enumi})}
\setcounter{enumi}{9}
\tightlist
\item
  我们在 \(\mathbf{N}^{*}\) 上定义 von Mangoldt 函数如下：\(\Lambda(n)=0\)（若 n 不是形如 \(p^{k}\) 的数，其中 p 是素数且 \(k\geqslant1\)），而 \(\Lambda(p^{k})=\log(p)\)。
\end{enumerate}

我们用 \(\zeta\) 表示 \(\mathbf{C} - \{1\}\) 上的全纯函数，它到满足如下条件的 \(s \in \mathbf{C}\) （\(\mathcal{R}(s) > 1\)）组成的集合上的限制与 Riemann zeta 函数重合（II，第 266 页，练习 16）。

对 \(x > 0\)，我们置

\[
\psi (x) = \sum_ {1 \leqslant n \leqslant x} \Lambda (n).
\]

\begin{enumerate}
\def\labelenumi{\alph{enumi})}
\tightlist
\item
  证明：对所有 \(s \in \mathbf{C}\)（满足 \(\mathcal{R}(s) > 1\)），有
\end{enumerate}

\[
\sum_ {n \geqslant 1} \Lambda (n) n ^ {- s} = - \frac {\zeta^ {\prime} (s)}{\zeta (s)}.
\]

\begin{enumerate}
\def\labelenumi{\alph{enumi})}
\setcounter{enumi}{1}
\tightlist
\item
  我们有
\end{enumerate}

\[
\lim _ {x \to + \infty} \frac {\psi (x)}{x} = 1.
\]

（利用 II，第 266 页，练习 16 以及练习 9）。

\begin{enumerate}
\def\labelenumi{\alph{enumi})}
\setcounter{enumi}{2}
\tightlist
\item
  设 \(\pi(x)\) 为 \(\leqslant x\) 的素数的个数。证明
\end{enumerate}

\[
\pi (x) \sim \frac {x}{\log (x)}
\]

（当 \(x \to +\infty\)）（“素数定理”，归于 Hadamard 与 de la Vallée Poussin）。

\begin{enumerate}
\def\labelenumi{\arabic{enumi})}
\setcounter{enumi}{10}
\tightlist
\item
  我们如下定义 Möbius 函数 \(\mu\)（在 \(\mathbf{N}^*\) 上）：\(\mu(n) = (-1)^k\)（若 \(n\) 是 \(k\) 个两两不同的素因子之积），否则 \(\mu(n) = 0\)。
\end{enumerate}

对 \(x > 0\)，我们置

\[
\mathrm{M} (x) = \sum_ {n \leqslant x} \mu (n), \quad f (x) = \frac {\mathrm{M} (x)}{x}.
\]

\begin{enumerate}
\def\labelenumi{\alph{enumi})}
\tightlist
\item
  证明：对 \(\mathcal{R}(s) > 1\)，有
\end{enumerate}

\[
\sum_ {n \geqslant 1} \mu (n) n ^ {- s} = \frac {1}{\zeta (s)}.
\]

\begin{enumerate}
\def\labelenumi{\alph{enumi})}
\setcounter{enumi}{1}
\item
  证明 f 是有界的，且 \(f(y) - f(x) = \mathrm{O}((y - x)/x)\)（对 \(y \geqslant x \geqslant 1\)），从而 f 在群 \(R_{+}^{*}\) 上当 x 趋于 \(+\infty\) 时是缓变的。
\item
  设 \(a, b \in R_{+}^{*}\)。对 \(x \textgreater{} 0\)，设 {[}x{]} 为 x 的整数部分。我们置 \(g_{0}(t) = [t^{-1}]\)（对 \(t \in R_{+}^{*}\)），并置
\end{enumerate}

\[
g (t) = 2 g _ {0} (t) - a g _ {0} (a t) - b g _ {0} (b t).
\]

函数 g 在 \(R_{+}^{*}\) 上关于 Lebesgue 测度可积。若 \(s \in C\) 且 \(\mathcal{R}(s) > 0\)，则

\[
\int_ {0} ^ {+ \infty} g (t) t ^ {s} d t = (2 - a ^ {- s} - b ^ {- s}) \frac {\zeta (1 + s)}{1 + s}.
\]

\(d)\) 我们有 \(\int_{0}^{+\infty} g(t)t^{ix} dt \neq 0\)（对所有 \(x \in \mathbf{R}\)）。（利用 II，第 266 页，练习 16 的 \(j\) )。）

\begin{enumerate}
\def\labelenumi{\alph{enumi})}
\setcounter{enumi}{4}
\tightlist
\item
  证明 \(\int_0^{+\infty} g(t / x)f(t)dt = o(x)\)（当 \(x \to +\infty\)）。推出 \(\mathrm{M}(x) = o(x)\)（当 \(x \to +\infty\)）。
\end{enumerate}

（可以初等地证明这个结果与素数定理等价，参看练习 10；例如参看 H. IWANIEC and E. KOWALSKI, Analytic number theory, A.M.S Coll. Publ. 53 (2004), 第 31--32 页。）

\begin{enumerate}
\def\labelenumi{\arabic{enumi})}
\setcounter{enumi}{11}
\item
  设 J 为满足如下条件的 \(f \in \mathcal{S}(\mathbf{R}^3)\) 组成的集合：\(\mathcal{F}(f)\) 在球面 \(\mathbf{S}_2\) 上为零。设 I 为满足如下条件的 \(f \in \mathrm{J}\) 组成的集合：\((\partial/\partial y_1)(\mathcal{F}(f))\) 在 \(\mathbf{S}_2\) 上为零。设 \(\overline{\mathrm{I}}\) 与 \(\overline{\mathrm{J}}\) 为 I 与 J 在 \(\mathrm{L}^1(\mathbf{R}^3)\) 中的闭包。则我们有 \(\mathrm{V}(\overline{\mathrm{I}}) = \mathrm{V}(\overline{\mathrm{J}}) = \mathbf{S}_2\)，但 \(\overline{\mathrm{I}} \neq \overline{\mathrm{J}}\)。（设 \(\mu\) 为 \(\mathbf{S}_2\) 上质量为 1、在 \(\mathbf{R}^3\) 的正交群下不变的正测度；利用 II，第 269 页，练习 20 的 c)，证明 \(f(t) = t_1\mathcal{F}(\mu)(t)\) 是 \(\mathrm{L}^\infty(\mathbf{R}^3)\) 中与 I 正交但不与 J 正交的元素。）
\item
  我们称 \(\widehat{\mathbf{G}}\) 中的 C 集是指 \(\widehat{\mathbf{G}}\) 的具有下列性质的闭子集 E：若 \(f\in \mathrm{L}^{1}(\mathrm{G})\)，若 \(\mathcal{F}(f)\) 在 E 上为零，且若 \(\varepsilon >0\)，则存在函数 \(g\in \mathrm{L}^{1}(\mathrm{G})\)，使得 \(\| f - f*g\| \leqslant \varepsilon\) 且 \(\operatorname {Supp}(\mathcal{F}(g))\) 是紧的并与 E 不相交。
\end{enumerate}

\begin{enumerate}
\def\labelenumi{\alph{enumi})}
\item
  若 E 是 \(\widehat{G}\) 中的 C 集，则恰存在 \(\mathrm{L}^1 (\mathrm{G})\) 的唯一闭理想 I 使得 \(\mathrm{V(I)} = \mathrm{E}\)。
\item
  \(\widehat{\mathrm{G}}\) 的每个单点子集都是 C 集。
\item
  两个 C 集的并是 C 集。
\item
  由 C 集经平移得到的每个集合都是 C 集。
\item
  设 H 为 \(\widehat{\mathbf{G}}\) 的闭子群，E 为 H 的闭子集，E' 为 E 相对于 H 的边界。若 E' 是相对于 \(\widehat{\mathbf{G}}\) 的 C 集，则 E 是关于 \(\widehat{\mathbf{G}}\) 的 C 集。（利用 II，第 270 页，练习 21。）
\item
  在 \(R^{n}\) 中，有限个闭半空间的交 E 是 C 集。（对 E 生成的仿射子空间的维数用归纳法，利用 b)、c)、d)、e)。
\end{enumerate}

\begin{enumerate}
\def\labelenumi{\arabic{enumi})}
\setcounter{enumi}{13}
\tightlist
\item
  设 E 为 \(R^{n}\) 的闭子集。假设存在 E 的一个内点 p，使得每条过 p 的直线与 E 的边界至多交于两点。则恰存在 \(L^{1}(\mathbf{R}^{n})\) 的唯一闭理想 I 使得 \(\mathcal{F}(f)\) 在 E 上为零。（把 f 看作 f 经以 p 为中心的位似所得到的函数在 \(L^{1}(\mathbf{R}^{n})\) 中的极限。）
\end{enumerate}

¶ 15) 我们用 μ 表示 G 的一个 Haar 测度，用 B 表示 e 在 \(\widehat{G}\) 中的邻域滤子。

\begin{enumerate}
\def\labelenumi{\alph{enumi})}
\tightlist
\item
  设 U 为 e 在 G 中的一个可积的开邻域。存在函数 \(a \in \mathrm{L}^{1}(\mathrm{G})\)，\(a \geqslant 0\)，积分为 1，在 U 之外为零，使得 \(\mathcal{F}(a) \in \mathrm{L}^{1}(\widehat{\mathrm{G}})\)，并且
\end{enumerate}

\[
\int_ {\mathrm{G}} a (x) ^ {2} d x \leqslant \frac {2}{\mu (\mathrm{U})}.
\]

（设 V 为 e 的紧邻域，使得 \(V \subset U\)，\(\mu(U) < \sqrt{2}\mu(V)\)；设 W 为 e 的紧对称邻域，使得 \(VW^{2} \subset U\)；取 \(a = \lambda f * g\)，其中 \(\lambda\) 是一个 \textgreater{} 0 的常数，f、g 是集合 VW 与 W 的特征函数。）

\begin{enumerate}
\def\labelenumi{\alph{enumi})}
\setcounter{enumi}{1}
\tightlist
\item
  若 \(f \in \mathrm{L}^{\infty}(\mathrm{G})\)，我们用 \(\mathrm{A}(f)\) 表示 \(\widehat{G}\) 中属于由 f 生成的、平移不变的弱闭向量子空间（\(\mathrm{L}^{\infty}(\mathrm{G})\) 的）的元素所成的集合。若 \(f, g \in \mathrm{L}^{\infty}(\mathrm{G})\)，则
\end{enumerate}

\[
\mathrm{A} (f g) \subset \overline {{\mathrm{A} (f) \mathrm{A} (g)}}.
\]

（设 U 为 \(e\) 的邻域。证明 \(fg\) 是 \(\widehat{G}\) 中属于 \(\mathrm{A}(f)\mathrm{UA}(g)\mathrm{U}\) 的元素的线性组合的弱极限。）

\begin{enumerate}
\def\labelenumi{\alph{enumi})}
\setcounter{enumi}{2}
\tightlist
\item
  设 \(g \in \mathrm{L}^{\infty}(\mathrm{G})\)，K 为 \(\widehat{G}\) 的紧子集，f 为 \(\mathrm{L}^{1}(\mathrm{G})\) 中的函数，使得 \(\mathcal{F}(f)\) 在 \(\mathrm{A} = \mathrm{A}(g)\) 上及在 K 之外为零。对每个 \(V \in B\)，设 \(\omega(V)\) 为 \(|\mathcal{F}(f)|\) 在 AV 上的上确界。若
\end{enumerate}

\[
\operatorname * {l i m i n f} _ {\mathfrak {B}} \frac {\omega (\mathrm{V}) ^ {2} \mu ((\mathrm{AV} - \mathrm{A}) \cap \mathrm{K})}{\mu (\mathrm{V})} = 0,
\]

则 \(\langle f,g\rangle=0\)。（设 \(\varepsilon>0\)。存在 G 的紧子集 H 使得 \(\int_{G-H}|f|dx\leqslant\varepsilon\)。于是存在 e 在 \(\widehat{G}\) 中的开邻域 U，使得 \(|\langle x,\widehat{x}\rangle-1|\leqslant\varepsilon\)（对 \(x\in H\)，\(\widehat{x}\in U\)）。把 a) 应用于 \(\widehat{G}\) 与 U，得到一个函数 \(a\in L^{1}(\widehat{G})\)。设 \(b=\mathcal{F}(a)\in L^{1}(G)\)。利用 b)，证明 \(\mathcal{F}(bg)\) 在 AU 之外为零。另一方面，

\[
\left| \int_ {\mathrm{G}} f g d x - \int_ {\mathrm{G}} f g b d x \right| \leqslant \varepsilon (\| f \| _ {1} + 2) \| g \| _ {\infty}
\]

并且，由于 \(f, gb \in \mathrm{L}^{2}(\mathrm{G})\)，

\[
\int_ {\mathrm{G}} f g b d x = \int_ {\widehat {\mathrm{G}}} \mathcal {F} (f) \mathcal {F} (g b) d \widehat {x} = \int_ {(\mathrm{AU-A}) \cap \mathrm{K}} \mathcal {F} (f) \mathcal {F} (g b) d \widehat {x}.
\]

把最后这个表达式的绝对值用下式界定：

\[
2 \| g \| _ {\infty} \frac {\omega (\mathrm{U}) \sqrt {\mu ((\mathrm{AU} - \mathrm{A}) \cap \mathrm{K})}}{\sqrt {\mu (\mathrm{U})}}
\]

即可完成证明。）

\(d)\) 设 I 为 \(\mathrm{L}^1 (\mathrm{G})\) 的闭理想，\(\mathrm{A} = \mathrm{V}(\mathrm{I})\)，\(f\) 为 \(\mathrm{L}^1 (\mathrm{G})\) 中的函数，使得 \(\mathcal{F}(f)\) 在 A 上为零，S 为 \(\mathcal{F}(f)\) 不为零的点所成的集合，\(\mathrm{S}'\) 为 S 的边界，B 为含于 \(\mathrm{S} \cap \mathrm{A}\) 的最大完满集。对 \(\mathrm{V} \in \mathfrak{B}\)，设 \(\omega (\mathrm{V})\) 为 \(|\mathcal{F}(f)|\) 在 AV 上的上确界。我们假设 B 的每个点在 \(\widehat{\mathrm{G}}\) 中有一个邻域 K，使得

\[
\operatorname * {l i m i n f} _ {\mathfrak {B}} \frac {\omega (\mathrm{V}) ^ {2} \mu ((\mathrm{AV-A}) \cap \mathrm{K})}{\mu (\mathrm{V})} = 0.
\]

则 \(f \in I\)。（利用 \(L^1(G)\) 满足 Ditkin 条件这一事实，并按 I，第 94 页，命题 5 的方式论证；设 \(G' = \widehat{G} \cup \{\infty\}\) 为 \(\widehat{G}\) 的 Alexandroff 紧化，并设 N 为满足如下条件的 \(\chi \in G'\) 所成的集合：\(\mathcal{F}(f)\) 不属于 \(\mathcal{F}(I)\)（在 \(\chi\) 的一个邻域中）。首先证明 \(N \subset B \cup \{\infty\}\)。然后，利用 \(c\) )，证明 \(N \subset \{\infty\}\)。最后，证明 \(N = \varnothing\)。）

\begin{enumerate}
\def\labelenumi{\alph{enumi})}
\setcounter{enumi}{4}
\tightlist
\item
  设 I 为 \(\mathrm{L}^{1}(\mathbf{R}^{n})\) 的闭理想，A = V(I)，f 为 \(\mathrm{L}^{1}(\mathbf{R}^{n})\) 中的函数，使得 \(\mathcal{F}(f)\) 在 A 上为零。对 h \textgreater{} 0，设 \(A_{h}\) 为 \(R^{n}\) 中在 A 之外且与 A 的距离小于 h 的点所成的集合，并设 \(\omega(h) = \sup_{x \in \mathrm{A}_{h}} |\mathcal{F}(f)(x)|\)。若
\end{enumerate}

\[
\liminf _ {h \to 0} \frac {\omega (h) ^ {2} \mu (\mathrm{A} _ {h})}{h ^ {n}} = 0
\]

则 \(f\in \mathrm{I}\)。（利用 \(d\) )。）

\(f)\) 设 \(f \in \mathrm{L}^{1}(\mathbf{R})\)，\(\alpha \in]0,1[\)，并设 \(\int_{\mathbf{R}} |y|^{\alpha}| f(y)| dy < +\infty\)。则存在常数 \(k\)，使得 \(|\mathcal{F}(f)(x+h)-\mathcal{F}(f)(x)| \leqslant kh^{\alpha}\)（对所有 \(x \in \mathbf{R}\) 以及每个 \(h \in \mathbf{R}\)）。

\(g)\) 设 I 为 \(\mathrm{L}^1 (\mathbf{R})\) 的闭理想，\(f\) 为 \(\mathrm{L}^1 (\mathbf{R})\) 中的函数，使得 \(\mathcal{F}(f)\) 在 \(\mathrm{V(I)}\) 上为零。若

\[
\int_ {\mathbf {R}} | y ^ {1 / 2} f (y) | d y <   + \infty ,
\]

则我们有 \(f \in I\)。（利用 e) 与 f)。）

\begin{enumerate}
\def\labelenumi{\arabic{enumi})}
\setcounter{enumi}{15}
\tightlist
\item
  设 \(f \in \mathrm{L}^{\infty}(\mathrm{G})\)。我们称 \(f\) 满足谱综合条件，如果它属于 \(\mathrm{Y}(\mathrm{A}(\mathrm{W}_f))\)，其中 \(\mathrm{W}_f\) 是 \(\mathrm{L}^{\infty}(\mathrm{G})\) 中包含 \(f\) 的所有平移不变弱闭子空间的交。
\end{enumerate}

\begin{enumerate}
\def\labelenumi{\alph{enumi})}
\item
  设 \(\mu \in \mathcal{M}^1 (\widehat{\mathrm{G}})\) 且 \(f = \overline{\mathcal{F}}_{\widehat{\mathrm{G}}}(\mu)\in \mathrm{L}^{\infty}(\mathrm{G})\)。则 \(\mathrm{A}(\mathrm{W}_f)\) 是 \(\mu\) 的支集。此外，\(f\) 满足谱综合条件。
\item
  若 G 是离散的，则每个函数 \(f \in \mathrm{L}^{2}(\mathrm{G})\) 都满足谱综合条件。
\end{enumerate}

\begin{enumerate}
\def\labelenumi{\arabic{enumi})}
\setcounter{enumi}{16}
\tightlist
\item
  对每个 \(p \geqslant 1\)，我们把 \(\mathrm{L}^p(\mathbf{R}_+)\) 与 \(\mathrm{L}^p(\mathbf{R})\) 的一个闭子空间等同。设 \(f \in \mathrm{L}^1(\mathbf{R}_+)\) 为由 \(f(x) = 2e^{-x}\)（对 \(x \in \mathbf{R}_+\)）定义的函数，并设 \(g \in \mathrm{L}^1(\mathbf{R})\) 为满足 \(g(x) = f(-x)\) 的函数。
\end{enumerate}

\begin{enumerate}
\def\labelenumi{\alph{enumi})}
\item
  证明 \(f + g = f * g\)。
\item
  设 A（相应地，B）为 \(\mathrm{L}^{1}(\mathbf{R})\) 中由 f（相应地，由 f 与 g）生成的子代数。证明 A 在 \(\mathrm{L}^{1}(\mathbf{R}_{+})\) 中稠密。（设 \(h \in \mathrm{L}^{\infty}(\mathbf{R}_{+})\) 满足：\(\langle h, f^{*n} \rangle = 0\)（对每个整数 \(n \geqslant 1\)），其中 \(f^{*n}\) 是 f 的 n 次迭次卷积；证明 Laplace 变换
\end{enumerate}

\[
\mathrm{F} (z) = \int_ {0} ^ {+ \infty} e ^ {- x z} h (x) d x,
\]

对所有 \(z \in C\)（满足 \(\mathcal{R}(z) > 0\)）都有定义的这个变换为零。）推出 B 在 \(\mathrm{L}^{1}(\mathbf{R})\) 中稠密。

\begin{enumerate}
\def\labelenumi{\alph{enumi})}
\setcounter{enumi}{2}
\item
  设 C 为 \(\mathrm{L}^{1}(\mathbf{R})\) 的包含 A 的真闭子代数。我们有 \(1 \in \mathrm{Sp}_{\mathrm{C}}(f)\)。（证明：否则，a) 以及 C 中的全纯函数演算将蕴含 \(g \in C\)。）
\item
  存在 C 的特征标 \(\chi\) 使得 \(\chi(f)=1\)；我们有
\end{enumerate}

\[
\chi (\varphi) = \int_ {0} ^ {+ \infty} e ^ {- x} \varphi (x) d x
\]

（对每个函数 \(\varphi\in\mathrm{L}^{1}(\mathbf{R}_{+})\)）。

\begin{enumerate}
\def\labelenumi{\alph{enumi})}
\setcounter{enumi}{4}
\item
  存在范数为 1 的测度 \(\mu \in \mathcal{M}^1 (\mathbf{R})\)，使得 \(\chi (\varphi) = \int \widehat{\varphi} d\mu\)（对每个函数 \(\varphi \in \mathrm{C}\)）。
\item
  \(\mu\) 的 Fourier 变换是函数 \(x \mapsto e^{-|x|}\)。
\item
  设 \(\varphi \in \mathrm{C}\)。对每个 \(y > 0\)，有
\end{enumerate}

\[
\int_ {\mathbf {R}} \varphi (x) (e ^ {- | x + y |} - e ^ {- | x | - y}) d x = 0.
\]

（利用关系式 \(\chi (\varphi *\psi) = \chi (\varphi)\chi (\psi)\)（对 \(\psi \in \mathrm{L}^1 (\mathbf{R}_+)\)）。）

\begin{enumerate}
\def\labelenumi{\alph{enumi})}
\setcounter{enumi}{7}
\tightlist
\item
  我们有 \(C = L^{1}(\mathbf{R}_{+})\)。（于是，子代数 \(L^{1}(\mathbf{R}_{+})\) 是 \(L^{1}(\mathbf{R})\) 的极大闭子代数。）
\end{enumerate}

\begin{enumerate}
\def\labelenumi{\arabic{enumi})}
\setcounter{enumi}{17}
\tightlist
\item
  设 G 为一个离散的有序交换拓扑群。我们把 G 的群运算写成加法，并用 0 表示它的单位元。我们用 \(G^{+}\) 表示 G 中满足 \(x \geqslant 0\) 的元素 x 所成的集合。
\end{enumerate}

\begin{enumerate}
\def\labelenumi{\alph{enumi})}
\item
  空间 \(L_{+}^{1}\)（由函数 \(f \in \mathrm{L}^{1}(\mathrm{G})\)（在 \(G^{+}\) 之外为零）组成）是 \(\mathrm{L}^{1}(\mathrm{G})\) 的闭子代数。
\item
  我们假设 G 不是 Archimedean 的（TG，V，第 16 页，§3 的练习 1）。代数 L \(_{+}^{1}\) 在 L \(^{1}\) (G) 中不是极大的。
\end{enumerate}

以下假设 G 是 Archimedean 的。设 A 为 \(\mathrm{L}^{1}(\mathrm{G})\) 的包含 \(L_{+}^{1}\) 的闭子代数。对每个 \(x \in G\)，设 \(\varphi_{x} \in \mathrm{L}^{1}(\mathrm{G})\) 为 \(\{x\}\) 的特征函数。它在 \(\mathrm{L}^{1}(\mathrm{G})\) 中可逆，其逆为 \(\varphi_{-x}\)。

\begin{enumerate}
\def\labelenumi{\alph{enumi})}
\setcounter{enumi}{2}
\tightlist
\item
  存在 \(y \in G_{+}\) 使得 \(\varphi_{y} \in A\) 且 \(\varphi_{-y} = \varphi_{y}^{-1} \notin A\)，并且存在 A 的特征标 \(\chi \neq 0\) 使得 \(\chi(y) = 0\)。
\end{enumerate}

\(d)\) 对 G 中每个 \(x > 0\)，我们有 \(\chi(x) = 0\)。（存在整数 \(k \geqslant 1\) 使得 \(-y + kx > 0\)。）

\begin{enumerate}
\def\labelenumi{\alph{enumi})}
\setcounter{enumi}{4}
\item
  存在范数为 1 的测度 \(\mu \in \mathcal{M}^1 (\mathrm{G})\)，使得 \(\chi\) 是 \(\mu\) 在代数 A（视为 \(\mathcal{C}(\mathrm{G})\) 的子空间）上的限制。（对 \(f\in \mathrm{A}\)，有 \(|\chi (f)|\leqslant \varrho (f)\)，而且 \(\varrho (f) = \| f\|_{\infty}\)。）
\item
  设 \(\widehat{\mu}\colon\mathrm{G}\to\mathbf{C}\) 为 \(\mu\) 的逆 Fourier 变换。我们有 \(\widehat{\mu}(x)=0\)（若 \(x>0\)）。
\item
  测度 \(\mu\) 是 G 上的计数测度。
\item
  子代数 \(L_{+}^{1}\) 在 \(L^{1}(G)\) 中是极大的。
\end{enumerate}

\begin{enumerate}
\def\labelenumi{\arabic{enumi})}
\setcounter{enumi}{18}
\tightlist
\item
  对 \(f \in \mathrm{L}^1(\mathbf{R})\)，我们用 \(\mathrm{A}_f\) 表示 \(\mathrm{L}^1(\mathbf{R})\) 中由 \(f\) 生成的闭子代数。
\end{enumerate}

\begin{enumerate}
\def\labelenumi{\alph{enumi})}
\tightlist
\item
  存在 \(f \in \mathrm{L}^{1}(\mathbf{R})\) 使得 \(A_{f}\) 是 \(\mathrm{L}^{1}(\mathbf{R})\) 的真极大子代数。
\end{enumerate}

在下文中，我们假设 \(f \in \mathrm{L}^{1}(\mathbf{R})\) 是使得 \(A_{f}\) 为 \(\mathrm{L}^{1}(\mathbf{R})\) 的真极大子代数的函数。

\begin{enumerate}
\def\labelenumi{\alph{enumi})}
\setcounter{enumi}{1}
\item
  我们记 \(S = \{0\} \cup \widehat{f}(\mathbf{R}) \subset \mathbf{C}\)。证明函数 \(z \mapsto \overline{z}\) 不是 S 上多项式函数的一致极限。（首先证明 C-S 不是连通的。）
\item
  \(\mathrm{L}^{1}(\mathbf{R})\) 中由 f 与 \(\widetilde{f}\) 生成的子代数在 \(\mathrm{L}^{1}(\mathbf{R})\) 中稠密。
\item
  f 的 Fourier 变换是单射的且在 R 上不取零值。（证明 S 不可能有两个有界的连通分支。）
\end{enumerate}

\backmatter
\specialchapter{Fourier 理论的公式表}

\specialsection{函数空间}

设 G 是一个局部紧交换群。我们回顾 G 上的主要函数空间以及它们与 Fourier 变换的关系。

左列定义 G 上的一个函数空间 E，并指出它是 \(L^{1}(G)\) 的子空间、\(L^{2}(G)\) 的子空间，还是两者的子空间。右列对 E 在 Fourier 变换和余变换下的象指出类似的包含关系（后者依情形或是定义在 \(L^{1}(G)\) 上的那个，或是定义在 \(L^{2}(G)\) 上的那个）。记号 F 表示 Fourier 变换或 Fourier 余变换。

\(\mathcal{M}^1 (\mathrm{G})\)

\[
\mathcal {F} (\mathcal {M} ^ {1} (\mathrm{G})) \subset \mathcal {C} _ {b} (\widehat {\mathrm{G}})
\]

（有界测度的空间）

（命题 3，第 207 页）

\(\mathrm{L}^{1}(\mathrm{G})\)

\[
\mathcal {F} (\mathrm{L} ^ {1} (\mathrm{G})) \subset \mathcal {C} _ {0} (\widehat {\mathrm{G}})
\]

（可积函数）

\(\mathrm{L}^2 (\mathrm{G})\)

\[
\mathscr {F} (\mathrm{L} ^ {2} (\mathrm{G})) = \mathrm{L} ^ {2} (\widehat {\mathrm{G}})
\]

（平方可积函数）（定理 1，第 215 页）

\(\mathrm{B(G)}\subset \mathrm{L}^{1}(\mathrm{G})\) \(\mathcal{F}(\mathrm{B(G)}) = \mathrm{B}(\widehat{\mathrm{G}})\)

由满足如下条件的 \(f \in \mathrm{L}^1(\mathrm{G})\) 组成的空间（定理 3，第 222 页）

\[
\mathcal {F} _ {\mathrm{G}} (f) \in \mathrm{L} ^ {1} (\widehat {\mathrm{G}})
\]

\(\mathrm{A(G)}\subset \mathrm{L}^1 (\mathrm{G})\cap \mathrm{L}^2 (\mathrm{G})\cap \mathrm{B(G)}\) \(\mathcal{F}(\mathrm{A(G)})\subset \mathcal{C}_0(\widehat{\mathrm{G}})\cap \mathrm{L}^1 (\widehat{\mathrm{G}})\cap \mathrm{L}^2 (\widehat{\mathrm{G}})\) （由 \(f*g\) 张成的空间，其中 \(f\)、\(g\in \mathrm{L}^{1}(\mathrm{G})\cap \mathrm{L}^{2}(\mathrm{G}))\)）（命题 8 的推论，第 211 页，引理 5，第 215 页，命题 11，第 217 页）

\specialsection{\texorpdfstring{\(\mathbf{R}^n\) 中的公式表}{R\^{}n 中的公式表}}

我们在此回顾关于 \(R^{n}\)（赋以 Euclid 范数）上的 Fourier 变换的最重要的公式。我们把 \(R^{n}\) 与它的对偶等同（通过映射 \((x,y)\mapsto\exp(2i\pi x\cdot y)\)），其中 \(x\cdot y=\sum_{j}x_{j}y_{j}\) （对所有 \(x=(x_{j})\) 与 \(y=(y_{j})\)（属于 \(R^{n}\)））。

\[
\begin{array}{l l} f \in \mathrm{L} ^ {1} (\mathbf {R} ^ {n}) & \widehat {f} (y) = \int_ {\mathbf {R} ^ {n}} f (x) \exp (- 2 i \pi x \cdot y) d x \\ & (y \in \mathbf {R} ^ {n}; \text {Def. 3, p. 206}) \\ f \in \mathrm{L} ^ {1} (\mathbf {R} ^ {n}) \text {or L} ^ {2} (\mathbf {R} ^ {n}) & \widehat {f} _ {h} (y) = \exp (2 i \pi h \cdot y) \widehat {f} (y) \\ y \in \mathbf {R} ^ {n} & f _ {h} (x) = f (x + h) \\ & (\text {Eq. (11), p. 208}) \\ f \in \mathrm{L} ^ {1} (\mathbf {R} ^ {n}) \text {or L} ^ {2} (\mathbf {R} ^ {n}) & \widehat {g} _ {h} (y) = \widehat {f} (y - h) \\ h \in \mathbf {R} ^ {n} & g _ {h} (x) = f (x) \exp (2 i \pi h \cdot x) \\ & (\text {Eq. (12), p. 208}) \\ f \in \mathrm{L} ^ {1} (\mathbf {R} ^ {n}) \text {or L} ^ {2} (\mathbf {R} ^ {n}) & \widehat {f} \circ \sigma = \frac {1}{| \det (\sigma) |} \widehat {f} \circ {} ^ {t} \sigma^ {- 1} \\ \sigma \in \mathrm{GL} (n, \mathbf {R}) & (\text {Eq. (16), p. 208}) \\ f (x) = \exp (- Q (x)), & \widehat {f} (y) = \frac {1}{\det (\sigma)} \exp (- Q ^ {*} (y)), \\ Q (x) = \| \sigma (x) \| ^ {2} & Q ^ {*} (y) = \| ^ {t} \sigma^ {- 1} (y) \| ^ {2} \\ Q \text {positive definite quadratic form} & (\text {example, p. 239}) \\ f \in \mathrm{L} ^ {2} (\mathbf {R} ^ {n}) & \int_ {\mathbf {R} ^ {n}} | f (x) | ^ {2} d x = \int_ {\mathbf {R} ^ {n}} | \widehat {f} (y) | ^ {2} d y \\ & (\text {Th. 1, p. 215}) \\ f \in \mathrm{L} ^ {2} (\mathbf {R} ^ {n}) \text {such that} & f (x) = \int_ {\mathbf {R} ^ {n}} \widehat {f} (y) \exp (2 i \pi x \cdot y) d x \\ \widehat {f} \in \mathrm{L} ^ {1} (\mathbf {R} ^ {n}) & (\text {almost everywhere ; Prop. 12, p. 219}) \\ f \in A (\mathbf {R} ^ {n}) & f (x) = \int_ {\mathbf {R} ^ {n}} \widehat {f} (y) \exp (2 i \pi x \cdot y) d x \\ & (x \in \mathbf {R} ^ {n}; \text {Prop. 11 p. 217}) \end{array}
\]

\specialsection{\texorpdfstring{\((\mathbf{R}/a\mathbf{Z})^{n}\) 中的公式表}{(R/aZ)\^{}n 中的公式表}}

为了读者方便，我们对群 \((\mathbf{R}/a\mathbf{Z})^{n}\)（其中 a \textgreater{} 0，赋以总质量为 1 的 Haar 测度，记作 \(a^{-n} dx\)）叙述 Fourier 理论的主要公式。当 n = 1 时，这对应于周期为 a 的周期函数。\((\mathbf{R}/a\mathbf{Z})^{n}\) 的对偶群与群 \(Z^{n}\) 等同（通过映射 \((\mathbf{R}/a\mathbf{Z})^{n} \times \mathbf{Z}^{n} \to \mathbf{U}\)），该映射是由下列映射过渡到商而诱导的

\[
(x, y) \mapsto \exp \left(\frac {2 i \pi}{a} x \cdot y\right)
\]

（对所有 \(x \in R^{n}\) 以及每个 \(y \in Z^{n}\)）。\(Z^{n}\) 上的对偶测度于是就是计数测度。

\[
\begin{array}{l l} f \in \mathrm{L} ^ {1} ((\mathbf {R} / a \mathbf {Z}) ^ {n}) & \widehat {f} (y) = \frac {1}{a ^ {n}} \int_ {(\mathbf {R} / a \mathbf {Z}) ^ {n}} f (x) \exp \Bigl (- \frac {2 i \pi}{a} y \cdot x \Bigr) d x \\ & (y \in \mathbf {Z} ^ {n}; \text {Def. 3, p. 206}) \\ f \in \mathrm{L} ^ {1} ((\mathbf {R} / a \mathbf {Z}) ^ {n}) \text {or} & \widehat {f} _ {h} (y) = \exp (\frac {2 i \pi}{a} h \cdot y) \widehat {f} (y) \\ \mathrm{L} ^ {2} ((\mathbf {R} / a \mathbf {Z}) ^ {n}) & f _ {h} (x) = f (x + h) \\ h \in (\mathbf {R} / a \mathbf {Z}) ^ {n} & (\text {Eq. (11), p. 208}) \\ f \in \mathrm{L} ^ {1} ((\mathbf {R} / a \mathbf {Z}) ^ {n}) \text {or} & \widehat {g} _ {h} (y) = \widehat {f} (y - h) \\ \mathrm{L} ^ {2} ((\mathbf {R} / a \mathbf {Z}) ^ {n}) & g _ {h} (x) = f (x) \exp (\frac {2 i \pi}{a} h \cdot x) \\ h \in (\mathbf {R} / a \mathbf {Z}) ^ {n} & (\text {Eq. (12), p. 208}) \\ f \in \mathrm{L} ^ {2} ((\mathbf {R} / a \mathbf {Z}) ^ {n}) & \frac {1}{a ^ {n}} \int_ {(\mathbf {R} / a \mathbf {Z}) ^ {n}} | f (x) | ^ {2} d x = \sum_ {y \in \mathbf {Z} ^ {n}} | \widehat {f} (y) | ^ {2} \\ & (\text {Th. 1, p. 215}) \\ f \in \mathrm{L} ^ {2} ((\mathbf {R} / a \mathbf {Z}) ^ {n}) \text {such that} & f (x) = \sum_ {y \in \mathbf {Z} ^ {n}} \widehat {f} (y) \exp (\frac {2 i \pi}{a} y \cdot x) \\ \widehat {f} \in \mathrm{L} ^ {1} (\mathbf {Z} ^ {n}) & (\text {almost everywhere; Prop. 12, p. 219}) \\ f \in \mathrm{A} ((\mathbf {R} / a \mathbf {Z}) ^ {n}) & f (x) = \sum_ {y \in \mathbf {Z} ^ {n}} \widehat {f} (y) \exp (\frac {2 i \pi}{a} y \cdot x) \\ & (x \in (\mathbf {R} / a \mathbf {Z}) ^ {n}; \text {Prop. 11 p. 217}) \end{array}
\]

\specialsection{\texorpdfstring{\(\mathbf{R}^n\) 中的 Poisson 公式}{R\^{}n 中的 Poisson 公式}}

我们在此叙述 n = 1 这一特别重要的情形下的 Poisson 公式（II，第 230 页，命题 15 的推论）。

设 \(f\in \mathrm{L}^1 (\mathbf{R})\) 且 \(a > 0\)。假设

\[
\sum_ {h \in \mathbf {Z}} | \widehat {f} (a ^ {- 1} h) | <   + \infty ,
\]

\[
\sum_ {h \in \mathbf {Z}} | f (x + a h) | <   + \infty \text {   for all   } x \in \mathbf {R},
\]

\[
x \mapsto \sum_ {h \in \mathbf {Z}} f (x + a h) \text {   continuous   on   } \mathbf {R}.
\]

则，对所有 \(x \in \mathbf{R}\)，有

\[
\sum_ {h \in {\bf Z}} f (x + a h) = \frac {1}{a} \sum_ {k \in {\bf Z}} \widehat {f} \Bigl (\frac {k}{a} \Bigr) \exp (\frac {2 i \pi}{a} k x).
\]

\specialchapter{记号索引}

SpA(x) （谱）2 R(x,λ) （预解式）3 Ā （单位元的添入）4 Sp′A(x) （谱）4 X(A) （特征标）6 X(h) （特征标函子）6 G\_A(x) （Gelfand 变换）7 X′(A) （特征标）9 X′(h) （特征标函子）9 G′A(x) （Gelfand 变换）10 J(A) （本原理想）12 V(M) （J(A) 的闭集）12 Υ(T) （与 T 相伴的理想）13 cl(π) （π 的类）13 Ā （不可约表示的集合）14 γ （正则表示）16 δ （正则表示）16 B(X;K) （有界函数）17 C\_b(X;K) （有界连续函数）17 C\_0(X;K) （在无穷远处为零的连续函数）17 C(X;K) （连续函数）17 K(X;K) （具有紧支集的连续函数）17 M¹(G) （有界测度）19 ρ(x) （谱半径）21 ev\_x （在 x 处的取值）32 X′ （Alexandroff 紧化）33 Sp\^{}Λ\_A(x) （同时谱）42 O(U;F) （全纯函数）49 O(K;F) （芽）50 O(U;A) （全纯函数代数）50



\(\mathcal{O}(\mathrm{K};\mathrm{A})\) （芽代数）50 \(\Theta_{a}\) （全纯函数演算）51 \(f(a)\) （全纯函数演算的函数记号）72 \(j_{\mathrm{H}}\) （与 H 相伴的幂等元）82 \(\mathrm{R_H}(x,\lambda)\) （关于 \(\mathrm{A_H} = j_{\mathrm{H}}\mathrm{A}j_{\mathrm{H}}\) 的预解式）83 \(\mathcal{O}_{\mathbf{R}}(\mathrm{S})\) （芽的 \(\mathbf{R}\) -代数（满足 \(f = f^{*}\)））86 \(\mathrm{Sp}_{\mathrm{A(C)}}(x)\) （复谱）86 \(x^{*}\) （伴随）96 \(\mathrm{A}_h\) （Hermite 元）96 \(\mathcal{C}'(\mathrm{X})\) （在一点处为零的函数）103 \(f(x)\) （连续函数演算）110 \(\mathrm{A}_{+}\) （正元）115 \(x \geqslant y\) （\(\mathrm{A}_h\) 上的序关系）117 \(x^{+}\) （正部）117 \(x^{-}\) （负部）117 \(|x|\) （绝对值）117 \(x^{\alpha}\) （正元的幂）118 \(\sqrt{x}\) （正元的平方根）118\\
Stell(A) （包络星代数）124\\
Stell(G) （群的星代数）125 \(e_{\mathrm{H}}(u)\) （谱投影）129 \(\mathrm{E_H}(u)\) （谱子空间）129 \(\widetilde{\mathrm{E}}_{\mathrm{H}}(u)\) （谱子空间）129 \(\iota(u)\) （数值值域）135 \(m(u)\) （谱的下确界）138 \(\mathrm{M}(u)\) （谱的上确界）138 \(|u|\) （Hilbert 空间之间态射的绝对值）139 \((j,|u|)\) （极分解）140\\
P(X) （多项式的极限）146 \(\operatorname{lsc}(g)\) （换位子的长度）155 \(\mathbf{C}[\mathrm{G}]\) （群代数）183 \(\operatorname{Stell}_r(\mathrm{G})\) （约化星代数）184 \(\operatorname{Lip}_b(\mathrm{X})\) （有界 Lipschitz 函数）196 \(\widetilde{f}(x)\) （\(\mathrm{L}^1 (\mathrm{G})\) 上的对合）199 \(\widehat{\mathbf{G}}\) （对偶群）201 \(\chi (\mu)\) （\(\mathcal{M}^1 (\mathrm{G})\) 的特征标）202 \(\mathrm{A}^{\perp}\) （正交）205 \(\widehat{\varphi}\) （对偶态射）205 \(\mathcal{F}_{\mathrm{G}}(\mu)\) （Fourier 变换）207 \(\overline{\mathcal{F}}_{\mathrm{G}}(\mu)\) （逆 Fourier 变换）207 \(\mathcal{F}_{\mathrm{G}}(f)\) （Fourier 变换）208 \(\overline{\mathcal{F}}_{\mathrm{G}}(f)\) （逆 Fourier 变换）208\\
A(G) （由 \(f*g\) 生成的子空间）210 \(\widehat{\mathrm{A}}(\mathrm{G})\) （A(G) 在 \(\mathcal{F}\) 下的象）211

B(G) （由 \(f \in L^{1}(G)\)（满足 \(\mathcal{F}(f) \in L^{1}(\widehat{G})\)）组成的子代数）222 T （环面）235 \(\mathcal{S}(\mathbf{R}^{n})\) （Schwartz 函数）238 V(Λ) （余体积）239 \(\mathcal{S}'(\mathbf{R}^{n})\) （缓增广义函数）239 A(W) （W 中的特征标）258 Y(F) （生成的 L∞(G) 的子空间）258 s(t) （符号函数）261 \(\mathcal{P}(X)\) （概率测度）261 ζ(s) （zeta 函数）267 L(s, χ) （Dirichlet L-函数）268 L(G) （G 在 R 中的表示）307 Λ(n) （von Mangoldt 函数）312

\specialchapter{术语索引}

\textbf{A}

伴随，96 单位元的添入，4, 105 Banach 代数，15 对合 Banach 代数，100 正则 Banach 代数，89 对合的，96 赋范的，15 赋范对合的，100 赋范乘积，16 对合子代数，97 幺的，1 星代数，102 群的，125 包络，124 商，122 约化，184 μ-遍历映射，190 真偏映射，33 自伴（--- 部分），96

\textbf{B}

Bernoulli（--- 移位），191 Bernstein（--- 不等式），279 Bohr-Pál（--- 定理），286 Shilov 边界，171

\textbf{C}

连续函数演算，110 全纯的，51

单变量的全纯函数演算，74 特征标，6, 9 Dirichlet，268 Legendre，288 幺的，6 酉的，201 Cauchy-Davenport（--- 定理），276 Chebotarev（--- 定理），276 Cohen 分解定理，185 幂等元定理，272 Ditkin 条件，92 谱综合条件，316 对称收敛，261 凸（多项式 --- 集），44 Fourier 余变换，207 Fourier 余变换，207 余体积，239 Weyl 一致分布判别法，277

\textbf{D}

Bernoulli 移位，191 极分解，140 星代数中的，119 谱分解，130 球堆积的密度，

Dirichlet 特征标 ---，268 --- L-函数，268 --- 定理，269 缓增广义函数，239 拓扑零因子，23 对偶群 ---，201 --- 测度，214 --- 态射，205 --- 格，238 Pontryagin 对偶，220 群之间的，225

\textbf{E}

正元，115 球堆积，295 预解集，2 多项式凸包，45 包络，70 Schrödinger 方程，269 模 1 的一致分布，277 遍历（μ--- 映射），190 特征标空间，8 指数，78

\textbf{F}

Fejér（--- 定理），242 Fekete（--- 引理），20 Fermat（--- 定理），288 忠实（--- 表示），11 概率测度的特征函数，280 共轭调和的，280 第一类 Bessel 函数，269 Möbius 函数，313 Schwartz 函数，238 von Mangoldt 函数，312 Hermite 函数，263 Dirichlet L-函数，268 缓变的，253 Lipschitz，196 概周期的，292 符号，261 theta，297 Riemann zeta，267

相伴微分形式，52 Hermite 线性型，98 正的，183 模形式，299 Hasse--Davenport 公式，288 Plancherel 公式，217 Poisson 公式，231 Fourier 反演公式，219 --- 的 Fourier 余变换，207 --- 反演公式，219 --- 级数，241 --- 变换，207 Fuglede--Putnam（定理），106

\textbf{G}

Gauss（和 ---），287, 299 Gauss（律 ---），281 Gelfand（--- 变换），7, 10 Gelfand--Mazur（--- 定理），26 全纯函数的芽，49 Gibbs（--- 现象），299 对偶群，201 单生的，306 顺从的，291 螺线管的，306 对偶中的群，225

\textbf{H}

Hadamard--de la Vallée Poussin（--- 的定理），313 Hardy（不确定性原理），274 Hardy--Littlewood（--- 的定理），309 Hasse--Davenport（--- 的公式），288 Hausdorff--Toeplitz（--- 的定理），136 Hermite 元 ---，96 线性型 ---，98 Hildebrandt（--- 的定理），189

\textbf{I}

本原理想，11 正则的，10

与元素相伴的幂等元，82 Ikehara（--- 的 Tauber 型定理），312 自同态的数值象，135 Bernstein 不等式，279 大筛法的，294 van der Corput 的，277 等周的，284 对合，96 半线性的，95 对合代数 ---，96 Banach 代数 ---，100 赋范代数 ---，100 不可约（表示 ---），11

\textbf{J}

Jacobi（--- 和），288 --- 的符号，299 Jacobson（--- 拓扑），13

\textbf{K}

Kaplansky（--- 的定理），184 Kolmogorov（--- 的定理），290 Kronecker（--- 的定理），277

\textbf{L}

Laplace（--- 变换），312 Legendre（--- 的特征标），288 Fekete 引理，20 Lévy --- 的定理，67 --- 的连续性定理，281 对数，78 Poisson 律，282 二次互反律的，297 Gauss 的，281 稳定的换位子长度，

\textbf{M}

计数测度，200 规范化的 Haar 测度，200 概率测度，261 对偶测度，214 态射

对合代数的，97 星代数的，102 表示的，11 幺的，1

\textbf{N}

Pisot 数，295 素数（--- 的定理），313 正规（元），96

\textbf{O}

Oka--Weil（--- 的定理），68 正交，205 正交投影算子，133

\textbf{P}

Paley-Wiener（--- 的定理），278\\
永久性，156\\
Gibbs 现象，299\\
Plancherel\\
--- 的公式，217\\
--- 的定理，215\\
完全（子代数 ---），5\\
Poisson\\
--- 的公式，231\\
--- 的律，282\\
多项式凸\\
集 ---，44\\
包络 ---，45\\
Pontryagin（--- 的对偶定理），220\\
正\\
元 ---，115\\
系统地 --- 元，182\\
本原（理想 ---），11\\
原理\\
不确定性，273, 275\\
Hardy 的不确定性，274\\
Euler 乘积，267\\
算术数列（Dirichlet 的 --- 的定理），269\\
长度为 3 的算术数列，271\\
谱投影算子，129\\
真（偏映射 ---），33

\textbf{Q}

拟幂零（元 ---），21

R Rademacher-Menshov（--- 的定理），287 谱半径，21 子集的，157 正则（理想 ---），10 正则 Banach 代数 ---，89 表示 ---，16 正交关系，233 代数的表示，11 忠实的，11 不可约的，11 正则的，16 等价的表示，11 格 --- 的余体积，239 对偶的，238 预解式，3 --- 的极点，131 Riemann（--- 的 zeta 函数），267 Roth（--- 的定理），271 Runge（--- 的定理），69, 151

\textbf{S}

Sahakyan（--- 的定理），285 Salem（--- 的定理），295 Schwartz（--- 的函数），238 半线性（对合 ---），95 Fourier 级数，241 Gauss 和，287, 299 Jacobi 和，288 Salié 和，288 对合子代数，97 完全的，5 独立子集，265 谱子空间，129 单参数子群，307 谱 --- 分解，130 --- 投影算子，129 --- 半径，21 --- 子空间，129 谱，2, 4 复的，86 自同态的，127

指数的，173 同时，42 奇异的，158 适应序列，52 Følner 的，291 拓扑群的正合序列，225 对偶，226 Jacobi 符号，299 拓扑生成系，16

\textbf{T}

Bohr-Pál 的定理，286 Cauchy-Davenport 的，276 Chebotarev 的，276 Lévy 连续性的，281 Pontryagin 对偶的，220 Cohen 分解的，185 Fejér 的，242 Fermat 的，288 Fuglede--Putnam 的，106 Gelfand--Mazur 的，26 Hausdorff--Toeplitz 的，136 Hildebrandt 的，189 Kaplansky 的，184 Kolmogorov 的，290 Kronecker 的，277 Dirichlet 算术数列的，269 Oka--Weil 的，68 P. Lévy 的，67 Paley-Wiener 的，278 Plancherel 的，215 Rademacher-Menshov 的，287 Roth 的，271 Runge 的，69, 151 Sahakyan 的，285 Salem 的，295 Wiener 的，38 Cohen 幂等元的，272 素数的，313 von Neumann 的遍历，190 概率论的基本极限，281 Hardy--Littlewood 的 Tauber 型，309 Ikehara 的 Tauber 型，312

Wiener 的 Tauber 型，253 Jacobson 拓扑，13 弱的，8, 10 Fourier 的变换，207 Gelfand 的，7 Hankel 的，270 Laplace 的，312 Fourier 变换，207

\textbf{U}

幺代数 ---，1 态射 ---，1 酉

特征标 ---，201 元 ---，96 逼近单位，120 递增的，120

V 特征值，128 van der Corput（--- 的不等式），277 von Neumann（--- 的遍历定理），190

W Weyl（--- 的判别法），277 Wiener --- 的定理，38 --- 的 Tauber 型定理，253
